% !TeX root = ../../Book3.tex
\chapter{Eisenstein 级数与模形式}
\label{b3:ch:23}
\BookChapterNumber{b3:ch:23}
% 本章摘要（不计入编号）
\begin{chapterabstract}
格的换基给出模群作用，格和给出带权变换的全纯函数。本章建立全模群的基本域、Eisenstein 级数和模形式，证明价公式并由此求出模形式环。判别式与 \(j\)-不变量把这些解析计算解释为复环面及其参数空间的几何。
\end{chapterabstract}

% 本章目标（不计入编号）
\begin{learninggoals}
能够推导模群作用与格比例的关系，检查尖点全纯性，计算 Eisenstein 级数的 Fourier 系数；能够在价公式中正确计入固定点的权重，用消失阶判断恒等式，并解释 \(j\) 对复环面的分类。
\end{learninggoals}

同一个格换一组基，椭圆函数代表的曲线并没有变化，但 \(\tau\)、\(g_2\)、\(g_3\) 的数值会按不同规律改变。哪些量只反映坐标，哪些量真正区别复结构，是本章要回答的问题。前章的格缩放、留数计算和三次曲线表示构成直接基础。Apostol 的《Modular Functions and Dirichlet Series in Number Theory》\cite{Apostol1990Modular}提供格和与模函数的系统叙述；本文价公式及其代数后果也可与 Zagier 的《Elliptic Modular Forms and Their Applications》\cite{Zagier2008Elliptic}相互核对。

% !TeX root = ../../Book3.tex
\section{模群}
\label{b3:sec:2301}

前章把一个格写成 \(\Lambda=\mathbb Z\omega_1+\mathbb Z\omega_2\)，并由它构造复环面与椭圆函数。若选取正向基，再除以 \(\omega_1\)，格就化为
\[
 \Lambda_\tau=\mathbb Z+\mathbb Z\tau,
 \quad \tau=\frac{\omega_2}{\omega_1}\in\mathfrak H,
 \quad \mathfrak H=\{\tau\in\mathbb C:\operatorname{Im}\tau>0\}.
\]
这里的上半平面与前文的 \(\mathbb H\) 是同一个区域；本章用 \(\mathfrak H\) 突出它的模参数角色。参数不是唯一的：改变同一个格的基，就会改变 \(\tau\)。模群先回答哪些参数表示同一比例格，再为后面的函数变换律提供准确的换基因子。

\subsecref{b3:subsec:110404}已从覆盖的角度讨论模群。本节恢复所需的矩阵与基本域事实，并进一步确定固定点和尖点的局部坐标。特别要记住，模群含有固定点，商映射不能直接套用自由群作用的覆盖结论。

\subsection{\texorpdfstring{\(\mathrm{SL}_2(\mathbb Z)\)}{SL2(Z)}}
\label{b3:subsec:230101}

\Needspace{6\baselineskip}
\begin{definition}\label{b3:def:23-modular-group}
行列式为一的二阶整数矩阵组成群
\[
 \operatorname{SL}_2(\mathbb Z)
 =\left\{\begin{pmatrix}a&b\\c&d\end{pmatrix}:
       a,b,c,d\in\mathbb Z,\ ad-bc=1\right\},
\]
群运算为矩阵乘法。本章在记录带权变换律时使用这个矩阵群；在只讨论其分式线性作用时，模群指它的有效作用群
\(\operatorname{PSL}_2(\mathbb Z)=\operatorname{SL}_2(\mathbb Z)/\{\pm I\}\)，与\subsecref{b3:subsec:110404}的约定一致。
\end{definition}
\termidx[moqun]{模群}\foreignidx{modular group}

逆矩阵为 \(\left(\begin{smallmatrix}d&-b\\-c&a\end{smallmatrix}\right)\)，所以整数条件在取逆后仍保持。行列式一还给一个将反复使用的事实：底行 \((c,d)\) 的两个整数互素。反过来，若 \(c,d\) 互素，Bézout 恒等式便给整数 \(a,b\) 使 \(ad-bc=1\)，从而把这行补成一个群元素。

记
\[
 T=\begin{pmatrix}1&1\\0&1\end{pmatrix},
 \quad S=\begin{pmatrix}0&-1\\1&0\end{pmatrix}.
\]
\cref{b3:prop:11-modular-generators}的 Euclid 算法也适用于这里的矩阵版本。若第一列为 \((a,c)^{\mathsf T}\)、\(c\ne0\)，作带余除法 \(a=qc+r\)、\(|r|<|c|\)，左乘 \(T^{-q}\) 和 \(S\) 将它依次变成
\[
 (a,c)^{\mathsf T}\longmapsto(r,c)^{\mathsf T}
 \longmapsto(-c,r)^{\mathsf T}.
\]
下左元素的绝对值严格下降，有限步后为零。所得矩阵的两个对角元素只能同为 \(1\) 或同为 \(-1\)，故为某个 \(\pm T^n\)。由于 \(-I=S^2\)，逆向读取操作就把原矩阵写成 \(S,T\) 及其逆的有限乘积。这个生成结论以后能把“对所有整数矩阵检查变换律”缩成对两个基本变换的检查。

\subsection{\texorpdfstring{\(\mathrm{PSL}_2(\mathbb Z)\)}{PSL2(Z)}}
\label{b3:subsec:230102}

矩阵 \(\gamma=\left(\begin{smallmatrix}a&b\\c&d\end{smallmatrix}\right)\) 决定分式线性映射
\[
 \gamma\tau=\frac{a\tau+b}{c\tau+d}.
\]
同时改变四个系数的符号不改变映射，所以 \(\gamma\) 和 \(-\gamma\) 作用相同。这个重复恰好是全部重复：若 \(\gamma\tau=\tau\) 在上半平面恒成立，则
\[
 c\tau^2+(d-a)\tau-b=0
\]
是一个恒为零的多项式，故 \(c=b=0\)、\(a=d\)；再由行列式条件得 \(a=d=\pm1\)。两个矩阵作用相同，便可对它们的商使用这个结论，说明它们只差整体负号。

因而 \(\operatorname{PSL}_2(\mathbb Z)\) 的每个元素是一个二元集合 \(\{\gamma,-\gamma\}\)，且它在上半平面上的作用忠实。这里取的是矩阵群的商；稍后还要对上半平面中的点按轨道取商，两种操作的对象不同。

矩阵等式为
\[
 S^2=-I,\quad (ST)^3=-I,
\]
所以在射影群中，\(S\) 的阶为二，\(ST\) 的阶为三。乘积 \(ST\) 表示先作 \(T\)、后作 \(S\)，其分式为 \(-1/(\tau+1)\)。在带权函数的变换律中仍保留矩阵 \(-I\)：虽然它不移动 \(\tau\)，它的权因子未必是一，这将直接排除本章的奇数权形式。

\subsection{在 \texorpdfstring{\(\mathfrak H\)}{上半平面} 上的作用}
\label{b3:subsec:230103}

先核对作用确实留在所需区域内。若 \(c\ne0\)，分母的零点 \(-d/c\) 在实轴上；若 \(c=0\)，则 \(d=\pm1\)。因此分母在 \(\mathfrak H\) 不为零。将分式与其共轭相减，并使用 \(ad-bc=1\)，得到
\begin{equation}\label{b3:eq:23-imaginary-transformation}
 \operatorname{Im}(\gamma\tau)
 =\frac{\operatorname{Im}\tau}{|c\tau+d|^2},
 \quad \gamma'(\tau)=\frac1{(c\tau+d)^2}.
\end{equation}
前一式说明虚部仍为正，后一式说明它处处局部可逆；由整数逆矩阵给出的逆映射保证它是上半平面的全纯自同构。

现在计算换基与参数变换之间的对应。前章\cref{b3:prop:22-lattice-bases}证明，同一格的两组正向基恰由整数行列式一矩阵联系。为与上述分式的系数顺序一致，把新基写成
\begin{equation}\label{b3:eq:23-basis-change}
 \eta_1=d\omega_1+c\omega_2,
 \quad \eta_2=b\omega_1+a\omega_2.
\end{equation}
相应换基矩阵是 \(\left(\begin{smallmatrix}d&b\\c&a\end{smallmatrix}\right)\)，其行列式仍为一。于是
\[
 \frac{\eta_2}{\eta_1}
 =\frac{a\tau+b}{c\tau+d},
 \quad \eta_1=\omega_1(c\tau+d).
\]
后一个尺度因子与前一个参数变换同样重要：模形式的权将记录它怎样作用于格上的量。

\Needspace{6\baselineskip}
\begin{proposition}[参数轨道与比例格]\label{b3:prop:23-lattice-orbits}
设 \(\tau,\sigma\in\mathfrak H\)，\(\Lambda_\tau=\mathbb Z+\mathbb Z\tau\)、\(\Lambda_\sigma=\mathbb Z+\mathbb Z\sigma\)。存在 \(\lambda\in\mathbb C\setminus\{0\}\) 使 \(\Lambda_\sigma=\lambda\Lambda_\tau\)，当且仅当存在
\(\gamma=\left(\begin{smallmatrix}a&b\\c&d\end{smallmatrix}\right)\in\operatorname{SL}_2(\mathbb Z)\) 使 \(\sigma=\gamma\tau\)。对每个这样的 \(\gamma\)，都有
\begin{equation}\label{b3:eq:23-lattice-transformation}
 \Lambda_{\gamma\tau}=(c\tau+d)^{-1}\Lambda_\tau.
\end{equation}
\end{proposition}
\begin{proof}
整数对 \((d,c)\)、\((b,a)\) 构成行列式一的矩阵，所以 \(c\tau+d\)、\(a\tau+b\) 是 \(\Lambda_\tau\) 的一组基。将这组基同时除以 \(c\tau+d\)，所得格由 \(1\) 和 \(\gamma\tau\) 生成，证明了公式及充分性。

反过来，若 \(\Lambda_\sigma=\lambda\Lambda_\tau\)，则 \((1,\sigma)\) 与 \((\lambda,\lambda\tau)\) 都是该格的正向基。第二组的有向面积为 \(|\lambda|^2\operatorname{Im}\tau>0\)，所以这里不能出现行列式负一的换基。由\cref{b3:prop:22-lattice-bases}，有整数 \(a,b,c,d\) 满足 \(ad-bc=1\) 及
\[
 1=\lambda(d+c\tau),\quad
 \sigma=\lambda(b+a\tau).
\]
相除即得 \(\sigma=\gamma\tau\)。
\end{proof}

例如 \(T\tau=\tau+1\) 只把第二个生成元改为 \(\tau+1\)，格的点集本身没有改变。\(S\tau=-1/\tau\) 则在交换两条基方向后重新把第一条基归一到一，产生因子 \(\tau^{-1}\)。这里分类的是相差非零复比例的格；它与复环面双全纯分类的关系将在\secref{b3:sec:2306}证明。

\subsection{基本域}
\label{b3:subsec:230104}

为在每条参数轨道中选代表，先用 \(T\) 把实部限制在宽一的竖条带内。若所得点还在单位圆内，\(S\) 会增大它的虚部。由此得到候选区域
\[
 \mathcal F=\{\tau\in\mathfrak H:
           |\operatorname{Re}\tau|\le\tfrac12,\ |\tau|\ge1\}.
\]
它包含边界，因此边界上的代表可以重复。以下陈述把唯一性准确限制在内部。

\Needspace{6\baselineskip}
\begin{proposition}\label{b3:prop:23-fundamental-domain}
设 \(\operatorname{SL}_2(\mathbb Z)\) 以上述分式线性变换作用于 \(\mathfrak H\)，令 \(T(\tau)=\tau+1\)、\(S(\tau)=-1/\tau\)，并令
\[
 \mathcal F=\{\tau\in\mathfrak H:
       |\operatorname{Re}\tau|\le1/2,\ |\tau|\ge1\}.
\]
每条轨道与 \(\mathcal F\) 相交；\(\mathcal F\) 内部的不同点不在同一轨道。\(\mathcal F\) 中的边界识别由两类配对生成：两条竖边由 \(T\) 配对，单位圆弧的两半由 \(S\) 配对。
\end{proposition}
\begin{proof}
固定 \(\tau\in\mathfrak H\)。在全部互素整数对 \((c,d)\) 上，\(|c\tau+d|\) 有正的最小值。具体说，候选 \((0,1)\) 已给上界一；在 \(|c\tau+d|\le1\) 时，先由 \(|c|\operatorname{Im}\tau\le1\) 限制 \(c\)，再由 \(|c\operatorname{Re}\tau+d|\le1\) 限制 \(d\)，所以只有有限个候选，而且没有一个值为零。把取得最小值的底行补成群矩阵，由\eqref{b3:eq:23-imaginary-transformation}，其像具有整条轨道中的最大虚部。再作整数平移使实部进入 \([-1/2,1/2]\)。若此时模小于一，作用 \(S\) 会严格增大虚部，矛盾。因此所得点属于 \(\mathcal F\)。

为比较两个代表，先证明对 \(z\in\mathcal F\) 和任意互素整数对 \((c,d)\)，有
\[
 |cz+d|\ge1.
\]
若 \(c=0\)，互素性给 \(|d|=1\)。若 \(|c|\ge2\)，由 \(\operatorname{Im}z\ge\sqrt3/2\) 得 \(|cz+d|\ge\sqrt3>1\)。剩下 \(|c|=1\)，改变整体符号后只须看 \(|z+d|\)。当 \(d=0\) 时用 \(|z|\ge1\)；当 \(|d|\ge1\) 时，
\[
 |z+d|^2=|z|^2+2d\operatorname{Re}z+d^2
 \ge1-|d|+d^2\ge1.
\]
若 \(z\) 在内部，\(c\ne0\) 的各情形都严格大于一。

若 \(z,w\in\mathcal F\)、\(w=\gamma z\)，分别对 \(\gamma\) 和 \(\gamma^{-1}\) 使用这个估计，得到两点虚部相等及 \(|cz+d|=1\)。在内部只可能 \(c=0\)，此时变换为整数平移；两个实部都在 \((-1/2,1/2)\)，平移量只能为零。

在边界，\(c=0\) 的非零平移只能把一条竖边送到另一条。若 \(|c|=1,d=0\)，取等要求 \(|z|=1\)，相应变换是 \(S\) 后接整数平移；除端点外，这给圆弧两半的配对。\(|c|=|d|=1\) 还能取等的点只有实部为 \(\pm1/2\) 的圆弧端点；它们的识别已由这两类配对生成。其余情形均不取等，所以没有遗漏的边界识别。
\end{proof}

为把这幅基本域图变成商空间上的局部坐标，还需要说明远处的群元素不会无限次挤入同一个小邻域。这是\cref{b3:prop:11-discrete-halfplane-action}的紧集有限相交性质；在当前整数群中可以直接看到它的原因。

\Needspace{6\baselineskip}
\begin{proposition}[模群作用的紧集有限性]\label{b3:prop:23-proper-action}
设 \(\Gamma=\operatorname{PSL}_2(\mathbb Z)\) 作用于 \(\mathfrak H\)。对任意紧集 \(K,L\subset\mathfrak H\)，满足 \(\gamma K\cap L\ne\varnothing\) 的 \(\gamma\in\Gamma\) 只有有限个。特别地，每点的稳定子 \(\Gamma_p=\{\gamma\in\Gamma:\gamma p=p\}\) 都是有限群。
\end{proposition}
\begin{proof}
只须讨论两个紧集均非空的情形。若 \(\tau\in K\)、\(\gamma\tau\in L\)，两点的虚部都有统一正下界和有限上界。虚部变换公式先给 \(|c\tau+d|\) 的上界，随后由 \(\operatorname{Im}\tau\) 的正下界限制整数 \(c\)，再由 \(\operatorname{Re}\tau\) 的有界性限制整数 \(d\)。因此底行只有有限种可能。

固定一个互素底行 \((c,d)\) 和一个上排解 \((a_0,b_0)\) 后，全部上排解恰为 \((a_0+nc,b_0+nd)\)、\(n\in\mathbb Z\)。对应映射是固定映射 \(\gamma_0\) 后接平移 \(n\)。紧集 \(\gamma_0K\) 与 \(L\) 的实部均有界，故能使两者相交的 \(n\) 只有有限个。取 \(K=L=\{p\}\) 即得稳定子有限。
\end{proof}

令 \(\rho=e^{2\pi\mathrm{i}/3}=(-1+\sqrt3\,\mathrm{i})/2\)。基本域中真正有非平凡稳定子的轨道只有 \(\mathrm{i}\) 与 \(\rho\)。对于前者，固定点方程给 \(a=d\)、\(b=-c\)、\(a^2+c^2=1\)，在射影群中只剩 \(I,S\)。对于后者，利用 \(\rho^2+\rho+1=0\)，固定点方程给
\[
 b=-c,\quad d=a+c,\quad a^2+ac+c^2=1.
\]
等价地 \((2a+c)^2+3c^2=4\)，故只能有 \(c=0,\pm1\)；逐个代入得六个矩阵，即 \(\pm I,\pm ST,\pm(ST)^2\)。除去正负重复后，稳定子阶为三。基本域中其他点由上一证明的取等分类可排除非恒等稳定子：竖边上的平移不固定点，圆弧上 \(S\) 只有固定点 \(\mathrm{i}\)，额外的端点情形正是 \(\rho\) 的轨道。

\Needspace{6\baselineskip}
\begin{proposition}[固定点附近的商坐标]\label{b3:prop:23-stabilizer-chart}
设 \(\Gamma=\operatorname{PSL}_2(\mathbb Z)\)、\(p\in\mathfrak H\)，令 \(e=|\Gamma_p|\)。若 \(p\) 分别等价于 \(\mathrm{i}\)、\(\rho=e^{2\pi\mathrm{i}/3}\)，则 \(e\) 分别为二、三；其余点有 \(e=1\)。取
\[
 u_p(\tau)=\frac{\tau-p}{\tau-\overline p}.
\]
存在 \(r>0\)，使 \(U=\{\tau:|u_p(\tau)|<r\}\) 中两点模群等价，当且仅当它们由 \(\Gamma_p\) 的元素联系。这个局部商由
\[
 w=u_p(\tau)^e
\]
同胚到圆盘 \(|w|<r^e\)。在 \(U\) 上对 \(\Gamma_p\) 不变的全纯函数，恰能唯一写成 \(w\) 的全纯函数；亚纯函数有同样的下降性质。
\end{proposition}
\begin{proof}
稳定子阶数已在基本域中算出；其他代表的稳定子由共轭得到，阶数不变。映射 \(u_p\) 把上半平面双全纯映到单位圆盘，并把 \(p\) 送到零。每个稳定子元素在这个坐标中是固定零的圆盘自同构，因而由 Schwarz 引理及其逆向应用成为旋转。一个 \(e\) 阶有限旋转群恰由全部 \(e\) 次单位根组成：取其中最小正转角，反复相减说明所有转角都是它的整数倍，再用群的阶确定该转角为 \(2\pi/e\)。

取 \(p\) 附近一个紧闭小盘。由\cref{b3:prop:23-proper-action}，只有有限多个群元素能使它与自身相交。对于其中不固定 \(p\) 的元素，缩小 \(p\) 的邻域便能使该邻域与它的像不交；有限次缩小后，再取一个更小的 \(u_p\) 坐标圆盘 \(U\)。它被稳定子保持，而所有其他群元素的像均与它不交。因此 \(U\) 内的等价关系正是 \(u\sim\zeta u\)、\(\zeta^e=1\)。

映射 \(u\mapsto u^e\) 的纤维恰是这些旋转轨道，且它是开映射，所以诱导局部商与圆盘之间的同胚。若全纯函数的 Taylor 级数为 \(\sum_{n\ge0}a_nu^n\)，旋转不变性迫使 \(a_n=0\) 除非 \(e\mid n\)。保留的项组成关于 \(w=u^e\) 的收敛幂级数。亚纯情形在零点附近用 Laurent 级数作相同论证；其他位置局部可逆，故同样下降。
\end{proof}

在 \(p=\mathrm{i}\) 时，直接代入可得 \(u_p(S\tau)=-u_p(\tau)\)；在 \(p=\rho\) 时，\(ST\) 在 \(u_p\) 坐标中的旋转因子为 \(e^{-2\pi\mathrm{i}/3}\)。局部商坐标分别是平方和立方。对于不变亚纯函数，上半平面上的整数零点阶数与商坐标中的阶数因此可能不同；价公式中的 \(1/2\)、\(1/3\) 将准确记录这件事。

\subsection{尖点与基本域约化}
\label{b3:subsec:230105}

模群除了在内部识别点，还识别一些边界方向。扩充有理轴 \(\mathbb Q\cup\{\infty\}\) 中的点称为\term[jiandian]{尖点}[cusp]。全模群的尖点只有一条轨道：既约分数 \(a/c\) 可以补成行列式一矩阵的第一列，该矩阵把 \(\infty\) 送到 \(a/c\)；反过来，整数分式线性变换把 \(\infty\) 送到有理数或 \(\infty\)。这里说的是有理边界点，不能把整个实轴都加进来当作同一种局部端。

固定 \(\infty\) 的射影群元素恰为整数平移。确实，固定无穷远要求 \(c=0\)，于是 \(a=d=\pm1\)，分式成为 \(\tau\mapsto\tau+n\)。因此自然尝试用
\[
 q=e^{2\pi\mathrm{i}\tau}
\]
消去这个平移重复。要确认它也是完整模群商的局部坐标，还须排除其他群元素在高处额外识别点。

\Needspace{6\baselineskip}
\begin{proposition}[尖点附近的穿孔圆盘]\label{b3:prop:23-cusp-chart}
设 \(\Gamma=\operatorname{PSL}_2(\mathbb Z)\)、\(Y>1\)，令 \(\mathfrak H_Y=\{\tau:\operatorname{Im}\tau>Y\}\)。若 \(\tau,\sigma\in\mathfrak H_Y\) 模群等价，则 \(\sigma-\tau\in\mathbb Z\)。因此 \(\mathfrak H_Y\) 在轨道商 \(\mathfrak H/\Gamma\) 中的像由
\[
 q=e^{2\pi\mathrm{i}\tau}
\]
同胚到穿孔圆盘 \(0<|q|<e^{-2\pi Y}\)，局部反函数由复对数给出。
\end{proposition}
\begin{proof}
若 \(\sigma=\gamma\tau\) 且 \(c\ne0\)，写 \(y=\operatorname{Im}\tau>Y\)，则
\[
 \operatorname{Im}\sigma
 =\frac{y}{|c\tau+d|^2}
 \le\frac1{c^2y}<\frac1Y<Y,
\]
与 \(\sigma\in\mathfrak H_Y\) 矛盾。所以只能 \(c=0\)，即整数平移。指数映射在 \(\mathfrak H_Y\) 上的像正是所述穿孔圆盘，两个原像的差为整数；它局部双全纯且是开映射，故诱导同胚。选一个小圆盘上的对数分支，\(\tau=(2\pi\mathrm{i})^{-1}\log q\) 就给局部逆；更换分支只会把 \(\tau\) 改变一个整数。
\end{proof}

在这个穿孔圆盘中加入 \(q=0\)，便把唯一的高端补成一个点。\secref{b3:sec:2306}将把它与刚才的内部坐标拼成紧 Riemann 曲面。“在尖点全纯”所要求的是关于 \(q\) 的函数在整个零点邻域全纯，而不是只沿某一条竖直线有有限极限。

\ProseParagraphBreak
基本域的证明也给出一个可执行的约化过程：先作整数平移，把实部送进 \([-1/2,1/2]\)；若模小于一，再作 \(S\)；重复这两步。每次真正使用 \(S\) 都严格增大虚部，而且不会无限重复。若初始点为 \(\tau_0\)、虚部为 \(y_0\)，每个途中点 \(\gamma\tau_0\) 的虚部至少为 \(y_0\)，于是
\[
 |c\tau_0+d|^2
 =\frac{y_0}{\operatorname{Im}(\gamma\tau_0)}\le1.
\]
这样的整数底行只有有限多个，所以途中可能出现的高度只有有限多个。严格增高的步骤只能有限次，最后便落入 \(\mathcal F\)。

例如 \(\tau=(1+2\mathrm{i})/10\) 已在竖条带内，却位于单位圆内。计算得
\[
 S\tau=-2+4\mathrm{i},\quad T^2S\tau=4\mathrm{i},
 \quad T^2S=\begin{pmatrix}2&-1\\1&0\end{pmatrix}.
\]
因此它与 \(4\mathrm{i}\) 表示同一比例格。换基公式给
\(\Lambda_{4\mathrm{i}}=\tau^{-1}\Lambda_\tau\)，而 \(\tau^{-1}=2-4\mathrm{i}\)，所以右边由 \(2-4\mathrm{i}\) 和一生成，确实等于 \(\mathbb Z+4\mathrm{i}\mathbb Z\)。约化改变的是格的归一化描述。内部点具有唯一的基本域代表；遇到边界时，按已证明的配对识别即可。

% !TeX root = ../../Book3.tex
\section{Eisenstein 级数}
\label{b3:sec:2302}

前章的格不变量随基变换产生确定的幂次因子。绝对收敛是把这种格对称性转成函数变换律的条件；权二处会失去它。

\subsection{格和}
\label{b3:subsec:230201}

对整数 \(k\ge2\)，定义\term[eisenstein-jishu]{Eisenstein 级数}[Eisenstein series]
\[
G_{2k}(\tau)\coloneqq\sum_{(m,n)\in\mathbb Z^2\setminus\{(0,0)\}}(m\tau+n)^{-2k}.
\]
这就是格 \(\Lambda_\tau\) 的 \(2k\) 次倒数和。它使用全部非零格向量；若只对互素 \((m,n)\) 求和，结果须乘以 \(\sum_{r\ge1}r^{-2k}\)，不能混用两种归一化。

\subsection{收敛}
\label{b3:subsec:230202}

在任意紧集 \(K\subset\mathfrak H\) 上，实线性变换 \((m,n)\mapsto m\tau+n\) 的最小奇异值有共同正下界，所以 \(|m\tau+n|\ge c_K\sqrt{m^2+n^2}\)。前章格壳估计给正规绝对收敛，故 \(G_{2k}\) 全纯。求导后的项含因子 \(m\)，但总衰减仍至少为 \((m^2+n^2)^{-k}\)，同一紧集估计允许任意固定次求导。

\ProseParagraphBreak

若 \(k=1\)，壳估计只给调和级数 \(\sum j^{-1}\)，绝对收敛失败。规定先对 \(n\) 再对 \(m\) 求和会得到有用函数，但它的变换律多出一项，不能把下面的重排论证照搬过去。

\subsection{\texorpdfstring{\(G_{2k}\)}{G2k}}
\label{b3:subsec:230203}

\begin{proposition}\label{b3:prop:23-eisenstein-transformation}
若 \(k\ge2\)、\(\gamma=\left(\begin{smallmatrix}a&b\\c&d\end{smallmatrix}\right)\in\mathrm{SL}_2(\mathbb Z)\)，则
\(G_{2k}(\gamma\tau)=(c\tau+d)^{2k}G_{2k}(\tau)\)。
\end{proposition}
\begin{proof}
每项分母为 \((m(a\tau+b)+n(c\tau+d))/(c\tau+d)\)。变换 \((m,n)\mapsto(ma+nc,mb+nd)\) 是 \(\mathbb Z^2\) 的双射，且保持非零对。绝对收敛允许依此重排，提出公共因子即得结果。
\end{proof}
令 \(\sigma_r(N)=\sum_{d\mid N}d^r\)。从 \(\pi\cot\pi z=-\pi\mathrm{i}(1+2\sum_{r\ge1}e^{2\pi\mathrm{i}rz})\) 在上半平面的展开出发，求 \(2k-1\) 次导数得到
\[
\sum_{n\in\mathbb Z}(z+n)^{-2k}
=\frac{(2\pi\mathrm{i})^{2k}}{(2k-1)!}\sum_{r\ge1}r^{2k-1}e^{2\pi\mathrm{i}rz}.
\]
先分出 \(m=0\)，将正负 \(m\) 配对，再按 \(N=mr\) 收集系数，得
\[
G_{2k}=2\zeta(2k)+\frac{2(2\pi\mathrm{i})^{2k}}{(2k-1)!}
\sum_{N\ge1}\sigma_{2k-1}(N)q^N.
\]
这里 \(\zeta(2k)\) 只表示收敛实级数。双和的绝对收敛由 \(|q|<1\) 和多项式系数增长保证，因而合并合法。

\subsection{与 \texorpdfstring{\(\wp\)}{\AnalysisPdfWeierstrassP} 的关系}
\label{b3:subsec:230204}

前章的记号现在为 \(g_2(\Lambda_\tau)=60G_4(\tau)\)、\(g_3(\Lambda_\tau)=140G_6(\tau)\)。格缩放律已经预示它们的权分别为四、六。相比之下，\(\wp(z;\Lambda_\tau)\) 还带一个曲线上的变量 \(z\)；改变 \(\tau\) 的基时必须同时缩放 \(z\)，不能将它误当成只依赖模参数的模形式。

% !TeX root = ../../Book3.tex
\section{模形式}
\label{b3:sec:2303}

Eisenstein 级数已经给出一种特殊的全纯函数：改变格基时，它不完全不变，而是乘上一个已知的幂次因子。同时，它在 \(\operatorname{Im}\tau\to\infty\) 的退化方向有有限的常数项。本节把这两项性质提炼为定义，并说明每一项的解析含义。尤其要分清三个层次：整数平移给周期性，整个模群给带权变换律，尖点条件则控制周期函数在无穷远的增长。

\subsection{模形式}
\label{b3:subsec:230301}

由\cref{b3:prop:23-eisenstein-transformation}，\(G_{2r}\) 在变换 \(\tau\mapsto(a\tau+b)/(c\tau+d)\) 下乘以 \((c\tau+d)^{2r}\)。这个因子来自格的尺度变化。若强求每个格和都在换基后数值相同，就会把已经构造出的有用对象排除。因而定义容许一个固定幂次的变化，但要求它对每个群元素都使用同一种规则。

另一方面，变换 \(T\tau=\tau+1\) 的因子为一，故所有这种函数都应当周期一。周期性使 \(f(\tau)\) 能用 \(q=e^{2\pi\mathrm{i}\tau}\) 表示：两个 \(\tau\) 给出同一个 \(q\)，当且仅当它们相差整数。由局部对数构成的函数先定义在穿孔圆盘上；要求它在 \(q=0\) 全纯，是一项额外条件。以下定义将这一要求直接写成一个圆盘全纯函数的存在性，避免把“尖点处全纯”留作未说明的边界用语。

\Needspace{6\baselineskip}
\begin{definition}\label{b3:def:23-modular-form}
设 \(k\in\mathbb Z\)，\(\mathfrak H=\{\tau\in\mathbb C:\operatorname{Im}\tau>0\}\)，\(\mathbb D=\{q\in\mathbb C:|q|<1\}\)。一个全纯函数 \(f:\mathfrak H\to\mathbb C\) 称为全模群 \(\operatorname{SL}_2(\mathbb Z)\) 上权 \(k\) 的\term[mo-xingshi]{模形式}[modular form]，是指它满足：
\begin{enumerate}
\item\label{b3:item:23-modular-transformation}对每个 \(\left(\begin{smallmatrix}a&b\\c&d\end{smallmatrix}\right)\in\operatorname{SL}_2(\mathbb Z)\) 和每个 \(\tau\in\mathfrak H\)，都有
\[
 f\!\left(\frac{a\tau+b}{c\tau+d}\right)
 =(c\tau+d)^k f(\tau);
\]
\item\label{b3:item:23-modular-cusp}存在全纯函数 \(F:\mathbb D\to\mathbb C\)，使
\[
 f(\tau)=F(e^{2\pi\mathrm{i}\tau})
 \quad(\tau\in\mathfrak H).
\]
\end{enumerate}
第二项称为在尖点 \(\infty\) 全纯。所有这样的函数构成复向量空间，记为 \(M_k\) 或 \(M_k(\operatorname{SL}_2(\mathbb Z))\)。本章不另加特征因子。
\end{definition}
\symidx{\(M_k\)：全模群上权 \(k\) 的模形式空间}

这里允许任意整数 \(k\) 是为了统一陈述。哪些权实际有非零形式，要由变换律和后面的零点计数判断；不应在定义时就把这些结论暗中假定。由于 \(c\tau+d\ne0\)，负整数幂也有单值的通常意义，没有复幂分支问题。

定义中的 \(F\) 必唯一，因为指数映射的像包含整个穿孔圆盘，两个可能的 \(F\) 在那里相等，故在零点也相等。空间 \(M_k\) 对复数线性组合封闭：变换律是线性的，圆盘上的全纯延拓也可作相同线性组合。前节的 Fourier 展开已经证明，对每个整数 \(r\ge2\)，都有 \(G_{2r}\in M_{2r}\)；它们不是靠定义猜出的例子，而是绝对收敛格和产生的具体成员。

\subsection{权}
\label{b3:subsec:230302}

定义中的整数 \(k\) 称为\term[mo-xingshi-quan]{权}[weight]。要理解它为何能与群乘法相容，记
\[
 J(\gamma,\tau)=c\tau+d,
 \quad \gamma=\begin{pmatrix}a&b\\c&d\end{pmatrix}.
\]
直接相乘两个矩阵给出
\begin{equation}\label{b3:eq:23-factor-cocycle}
 J(\gamma_1\gamma_2,\tau)
 =J(\gamma_1,\gamma_2\tau)J(\gamma_2,\tau).
\end{equation}
例如，若 \(\gamma_1\) 的底行为 \((c_1,d_1)\)，\(\gamma_2\) 的两行为 \((a_2,b_2)\)、\((c_2,d_2)\)，右边等于
\[
 c_1(a_2\tau+b_2)+d_1(c_2\tau+d_2),
\]
正是乘积矩阵的底行给出的线性式。因此依次作两次变换所积累的权因子，与一次作复合变换的因子一致。

为了把这个运算写得紧凑，对 \(\mathfrak H\) 上的全纯或亚纯函数 \(f\)、整数 \(k\) 和 \(\gamma\in\operatorname{SL}_2(\mathbb Z)\)，记
\[
 (f|_k\gamma)(\tau)
 =J(\gamma,\tau)^{-k}f(\gamma\tau).
\]
这个记号在变换后除去权因子，使带权变换律可以写成不变性条件。公式\eqref{b3:eq:23-factor-cocycle}给
\[
 (f|_k\gamma_1)|_k\gamma_2=f|_k(\gamma_1\gamma_2).
\]
所以变换律正是 \(f|_k\gamma=f\)。它也解释了为什么只核对生成元就足够：满足对 \(S,T\) 的不变性，便依次满足它们的逆和全部有限乘积的不变性。

\Needspace{6\baselineskip}
\begin{proposition}[生成元判别与权的运算]\label{b3:prop:23-weight-operations}
设 \(k\in\mathbb Z\)，\(f\) 为 \(\mathfrak H\) 上全纯函数。它满足全模群的权 \(k\) 变换律，当且仅当
\begin{equation}\label{b3:eq:23-generators-form-law}
 f(\tau+1)=f(\tau),\quad
 f(-1/\tau)=\tau^k f(\tau)
 \quad(\tau\in\mathfrak H).
\end{equation}
奇数权的 \(M_k\) 仅含零函数。若 \(k,\ell\in\mathbb Z\)、\(f\in M_k\)、\(g\in M_\ell\)，则 \(fg\in M_{k+\ell}\)；若 \(g\not\equiv0\)，商 \(f/g\) 在上半平面及尖点亚纯，并满足权 \(k-\ell\) 的变换律。
\end{proposition}
\begin{proof}
必要性是代入 \(S,T\)。充分性由上述带权运算的复合公式及 \(S,T\) 的生成性得到。对矩阵 \(-I\)，变换律给 \(f(\tau)=(-1)^kf(\tau)\)，所以 \(k\) 为奇数时只能 \(f\equiv0\)。

两个变换律相乘便使权相加；圆盘全纯表达式的乘积仍全纯，故乘积在尖点也满足要求。商的变换律同样由相除得到。非恒零全纯函数的零点孤立，局部写成一个整数次幂乘无零全纯函数，所以商在有限点亚纯；在尖点对两个圆盘 Taylor 级数作同样的局部分解，便知商至多有有限阶极点。
\end{proof}

同权形式可以相加，不同权的乘积则进入新的权空间。例如 \(G_4G_6\) 的权为十。若要同时保留所有权，就要把它们组成带权标记的分次空间；不能把不同权相加后仍称为原来某个固定权的形式。同权商的因子相消，得到不变的亚纯函数，这将在\secref{b3:sec:2306}成为模函数的来源。

\ProseParagraphBreak
权的另一个解释来自未归一化的格。对格不变量 \(G_{2r}(\Lambda)=\sum_{\omega\in\Lambda\setminus\{0\}}\omega^{-2r}\)，同时把全部格向量乘以 \(\lambda\) 会使它乘以 \(\lambda^{-2r}\)。下述命题说明一般变换律也是这样一种尺度齐次性的坐标表达。

\Needspace{6\baselineskip}
\begin{proposition}[权与格的尺度]\label{b3:prop:23-lattice-weight}
设 \(k\in\mathbb Z\)，\(f:\mathfrak H\to\mathbb C\) 满足全模群的权 \(k\) 变换律。对复平面中的格 \(\Lambda\)，选取任一正向基 \((\omega_1,\omega_2)\)，定义
\[
 \Phi_f(\Lambda)=\omega_1^{-k}
          f\!\left(\frac{\omega_2}{\omega_1}\right).
\]
此值不依赖正向基的选择，且对每个 \(\lambda\in\mathbb C\setminus\{0\}\) 有
\[
 \Phi_f(\lambda\Lambda)=\lambda^{-k}\Phi_f(\Lambda).
\]
反过来，若一个复值格函数 \(\Phi\) 满足上述尺度律，且 \(f(\tau)=\Phi(\mathbb Z+\mathbb Z\tau)\) 在 \(\mathfrak H\) 全纯，则 \(f\) 满足权 \(k\) 的变换律；它是否是模形式，还须检查尖点全纯性。
\end{proposition}
\begin{proof}
若另一组正向基为 \((\eta_1,\eta_2)\)，由\eqref{b3:eq:23-basis-change}，存在 \(\gamma\in\operatorname{SL}_2(\mathbb Z)\) 使
\[
 \eta_1=\omega_1(c\tau+d),\quad
 \frac{\eta_2}{\eta_1}=\gamma\tau,
 \quad \tau=\frac{\omega_2}{\omega_1}.
\]
因此
\[
 \eta_1^{-k}f(\eta_2/\eta_1)
 =\omega_1^{-k}(c\tau+d)^{-k}f(\gamma\tau)
 =\omega_1^{-k}f(\tau).
\]
乘以非零复数 \(\lambda\) 后，\((\lambda\omega_1,\lambda\omega_2)\) 仍正向，基的比值不变而前因子乘以 \(\lambda^{-k}\)，得到尺度律。

反过来，由\eqref{b3:eq:23-lattice-transformation}及尺度律，
\[
 f(\gamma\tau)
 =\Phi((c\tau+d)^{-1}\Lambda_\tau)
 =(c\tau+d)^k\Phi(\Lambda_\tau)
 =(c\tau+d)^kf(\tau).
\]
这个计算没有使用 \(\operatorname{Im}\tau\to\infty\) 时的任何估计，所以不会自动给出尖点条件。
\end{proof}

格的点集满足 \(-\Lambda=\Lambda\)，所以奇数齐次权也会迫使 \(\Phi_f=0\)，与矩阵 \(-I\) 的判别相同。另一方面，权并不表示 \(f\) 作为变量 \(\tau\) 的普通齐次次数。缩放格与把参数 \(\tau\) 乘一个复数是两种不同操作。

普通导数也不会自动把权提高二。设 \(k\in\mathbb Z\)，\(f\) 在 \(\mathfrak H\) 全纯并满足全模群权 \(k\) 的变换律。对这一变换律求导，使用 \((\gamma\tau)'=(c\tau+d)^{-2}\)，得到
\begin{equation}\label{b3:eq:23-ordinary-derivative-law}
 f'(\gamma\tau)
 =(c\tau+d)^{k+2}f'(\tau)
   +kc(c\tau+d)^{k+1}f(\tau).
\end{equation}
额外的一项来自权因子的导数。\subsecref{b3:subsec:230503}将构造修正项使它相消；这里先把普通导数的真实变换律保留下来。

\subsection{尖点模形式}
\label{b3:subsec:230303}

若模形式的尖点常数项为零，它在退化端比一般模形式衰减得更快。这个性质对线性组合稳定，因而给出一个自然子空间。

\Needspace{6\baselineskip}
\begin{definition}\label{b3:def:23-cusp-form}
设 \(k\in\mathbb Z\)，\(f\in M_k(\operatorname{SL}_2(\mathbb Z))\)，并令 \(F\) 为单位圆盘上的唯一全纯函数，满足 \(f(\tau)=F(e^{2\pi\mathrm{i}\tau})\)。若 \(F(0)=0\)，则称 \(f\) 为全模群上权 \(k\) 的\term[jiandian-mo-xingshi]{尖点模形式}[cusp form]。这类函数构成的复向量空间记为 \(S_k\) 或 \(S_k(\operatorname{SL}_2(\mathbb Z))\)。
\end{definition}
\symidx{\(S_k\)：全模群上权 \(k\) 的尖点模形式空间}

常数项映射
\[
 a_0:M_k\longrightarrow\mathbb C,\quad f\longmapsto F(0)
\]
是线性的，其核恰为 \(S_k\)。因此 \(M_k/S_k\) 的维数至多为一。若已找到常数项非零的权 \(k\) 形式，将它除以常数项就得到 \(a_0=1\) 的形式，从任何同权 \(f\) 减去适当倍数后便进入 \(S_k\)。这一简单操作将成为后面模形式环归纳证明的第一步。

尖点模形式乘以任何全纯模形式仍在尖点消失，所以 \(S_kM_\ell\subset S_{k+\ell}\)。不过，单独的周期性和衰减还不够：\(f(\tau)=e^{2\pi\mathrm{i}\tau}\) 周期一且在高处趋零，但它不是权零模形式。例如
\[
 f(2\mathrm{i})=e^{-4\pi},\quad
 f(S(2\mathrm{i}))=f(\mathrm{i}/2)=e^{-\pi},
\]
不满足权零的 \(S\) 不变性。三个条件必须一起核对，不能由某一端的良好行为代替群变换律。

这里检查一个尖点就够，是因为\subsecref{b3:subsec:230105}已证明全模群的有理尖点只有一条轨道。若用某个 \(\gamma\) 把 \(\infty\) 送到另一尖点，应先用 \(f|_k\gamma\) 表示该处函数；全模群的变换律使它仍等于 \(f\)，所以得到同一尖点条件。对有限指标子群可能出现多条尖点轨道，不能把本章的简化照搬过去。

\subsection{Fourier 展开}
\label{b3:subsec:230304}

现在证明前面使用的 \(q\) 表达式，并把尖点条件变成可以直接检查的增长条件。先暂时不要求 \(S\) 变换律，只保留周期性，这样就能看清 Fourier 展开究竟依赖哪一步。

\Needspace{6\baselineskip}
\begin{proposition}[周期函数的尖点展开]\label{b3:prop:23-q-expansion}
设 \(f:\mathfrak H\to\mathbb C\) 全纯，且 \(f(\tau+1)=f(\tau)\)。则存在唯一全纯函数
\(F:\{q:0<|q|<1\}\to\mathbb C\)，满足 \(f(\tau)=F(e^{2\pi\mathrm{i}\tau})\)。其 Laurent 展开给
\[
 f(\tau)=\sum_{n\in\mathbb Z}a_ne^{2\pi\mathrm{i}n\tau},
\]
在 \(\mathfrak H\) 的紧集上局部一致收敛，并且对任意整数 \(n\) 和任意 \(y>0\)，
\begin{equation}\label{b3:eq:23-fourier-coefficients}
 a_n=e^{2\pi ny}\int_0^1
        f(x+\mathrm{i}y)e^{-2\pi\mathrm{i}nx}\,\mathrm{d}x.
\end{equation}
以下三项等价：\(F\) 在 \(q=0\) 可去；存在 \(Y>0,C<\infty\) 使 \(|f(x+\mathrm{i}y)|\le C\) 对所有 \(x\in\mathbb R,y>Y\) 成立；所有 \(n<0\) 的 \(a_n\) 都为零。条件成立时，级数为 \(\sum_{n\ge0}a_nq^n\)，在整个 \(|q|<1\) 内收敛，且
\[
 \sup_{x\in\mathbb R}|f(x+\mathrm{i}y)-a_0|
 =O(e^{-2\pi y})\quad(y\to\infty).
\]
\end{proposition}
\begin{proof}
给定 \(0<|q|<1\)，任选一个满足 \(q=e^{2\pi\mathrm{i}\tau}\) 的 \(\tau\in\mathfrak H\)，定义 \(F(q)=f(\tau)\)。另一个选择只会将 \(\tau\) 改变整数，周期性保证值不变。在每个不含零的小圆盘内可选取全纯对数，故 \(F(q)=f((2\pi\mathrm{i})^{-1}\log q)\) 局部全纯。这既证明存在性，也由指数映射的满射性证明唯一性。

Laurent 定理在环域 \(0<|q|<1\) 给唯一展开 \(F(q)=\sum_{n\in\mathbb Z}a_nq^n\)，在每个紧子环域上一致收敛。\(\mathfrak H\) 中紧集的指数像是这种环域中的紧集，所以拉回也局部一致。系数公式为
\[
 a_n=\frac1{2\pi\mathrm{i}}
        \int_{|q|=r}F(q)q^{-n-1}\,\mathrm{d}q,
 \quad 0<r<1.
\]
取 \(r=e^{-2\pi y}\)，按 \(q=e^{2\pi\mathrm{i}(x+\mathrm{i}y)}\)、\(0\le x\le1\) 参数化这条正向圆周，便得到\eqref{b3:eq:23-fourier-coefficients}。这个推导也说明结果与 \(y\) 无关：不同高度只是计算同一个 Laurent 系数的不同积分圆周。

若 \(F\) 在零点可去，它在某个闭小圆盘上有界；因为 \(|q|=e^{-2\pi y}\)，对应地 \(f\) 在整个高半平面有界。反过来，高半平面的一致界转成 \(F\) 在某个穿孔小圆盘上的界，由可去奇点定理得到全纯延拓。Laurent 定理又说明，可去性恰等价于全部负次系数为零。

延拓后，\(F(q)=a_0+qB(q)\)，其中 \(B\) 在单位圆盘全纯。在一个固定闭小圆盘上 \(|B|\) 有界，故 \(|F(q)-a_0|\le C|q|\)。代入 \(q=e^{2\pi\mathrm{i}(x+\mathrm{i}y)}\) 就得到最后的一致估计。
\end{proof}

因此，对一个已经满足全模群变换律的全纯函数，验证尖点条件可以改为在 \(|\operatorname{Re}\tau|\le1/2\)、\(\operatorname{Im}\tau>Y\) 上取得统一界；周期性会把这个界延伸到整个高半平面。只沿一条竖线检查极限并不具有同样作用。例如
\[
 f(\tau)=\exp(-e^{-2\pi\mathrm{i}\tau})
\]
全纯且周期一，它对应 \(F(q)=e^{-1/q}\)。沿 \(\tau=\mathrm{i}y\) 有 \(f(\mathrm{i}y)=\exp(-e^{2\pi y})\to0\)，但沿 \(\tau=1/2+\mathrm{i}y\) 则有 \(f(\tau)=\exp(e^{2\pi y})\to\infty\)。一个方向的有限极限不能排除零点处的本性奇点。这个例子只用于检验周期函数的尖点行为，不声称它满足整个模群的变换律。

对模形式，\eqref{b3:eq:23-fourier-coefficients}就是 Fourier 系数的复分析来源。级数在每个固定高度上按 \(x\) 展开成周期函数，又在 \(q\) 圆盘中是一条 Taylor 级数。这里没有调用 \(L^2\) Fourier 级数的完备性：全纯性已经给出比一般实变量周期函数更强的局部一致收敛。特别地，模形式由全部系数唯一确定，而尖点模形式恰是 \(a_0=0\) 的那些形式。

\subsection{在尖点处消失}
\label{b3:subsec:230305}

尖点全纯只说明没有负次幂。若进一步需要计数零点，就应记录第一个非零幂次。设 \(f\in M_k\) 非恒零，写
\[
 f(\tau)=\sum_{n\ge0}a_nq^n,
 \quad q=e^{2\pi\mathrm{i}\tau}.
\]
其尖点消失阶定义为
\[
 v_\infty(f)=\min\{n\ge0:a_n\ne0\}.
\]
非恒零保证这个最小值存在；若所有系数为零，Taylor 唯一性会使 \(f\equiv0\)。可约定零函数的阶数为 \(+\infty\)，但后面涉及对数导数的公式只用于非恒零函数。尖点形式满足 \(v_\infty(f)\ge1\)，而一般非零模形式允许阶数为零。

价公式还要处理在尖点有极点的函数，所以将同一计算保留到 Laurent 展开具有有限主部的情形。下面的一致估计是围道顶边取极限时的直接依据。

\Needspace{6\baselineskip}
\begin{proposition}[尖点阶数与对数导数]\label{b3:prop:23-cusp-order-estimate}
设 \(Y_0\in\mathbb R\)，\(f\) 是半平面 \(\operatorname{Im}\tau>Y_0\) 上非恒零的周期一亚纯函数，其关于 \(q=e^{2\pi\mathrm{i}\tau}\) 的函数在 \(q=0\) 亚纯。设该 Laurent 展开的首项为 \(a_mq^m\)，其中 \(m\in\mathbb Z\)、\(a_m\ne0\)。则存在 \(Y>Y_0\) 和常数 \(C>0\)，使 \(f\) 在 \(\operatorname{Im}\tau>Y\) 没有零点及极点，而且对所有 \(x\in\mathbb R,y>Y\)，都有
\begin{align}
 \left|\frac{f(x+\mathrm{i}y)}{a_m e^{2\pi\mathrm{i}m(x+\mathrm{i}y)}}-1\right|
 &\le Ce^{-2\pi y},\label{b3:eq:23-cusp-leading-term}\\
 \left|\frac{f'(x+\mathrm{i}y)}{f(x+\mathrm{i}y)}-2\pi\mathrm{i}m\right|
 &\le Ce^{-2\pi y}.\label{b3:eq:23-cusp-log-derivative}
\end{align}
这里 \(m=v_\infty(f)\)，极点的阶数按负整数计。
\end{proposition}
\begin{proof}
尖点亚纯性给出分解 \(F(q)=q^mA(q)\)，其中 \(A\) 在零点某邻域全纯，且 \(A(0)=a_m\ne0\)。取足够小的 \(r>0\)，使 \(A\) 在 \(|q|\le r\) 无零点；再略缩小圆盘，\(A'\) 与 \(1/A\) 都有统一界。因此
\[
 A(q)=a_m+O(|q|),\quad
 \frac{A'(q)}{A(q)}=O(1)
\]
对整个小圆盘一致成立。令 \(e^{-2\pi Y}<r\)，\(q\ne0\) 及 \(A(q)\ne0\) 说明 \(f\) 在所述高半平面无零极点。第一式来自 \(F(q)/(a_mq^m)=A(q)/a_m\)。

对第二式，利用 \(q'(\tau)=2\pi\mathrm{i}q\) 直接计算
\[
 \frac{f'(\tau)}{f(\tau)}
 =2\pi\mathrm{i}\left(m+q\frac{A'(q)}{A(q)}\right).
\]
余项绝对值不超过常数乘 \(|q|=e^{-2\pi y}\)，与辐角 \(2\pi x\) 无关，因而估计对全部 \(x\) 一致。
\end{proof}

这比“每条竖线上趋于某值”更强：它同时控制整条水平线，因此可沿长度为一的顶边积分并让高度趋于无穷。对 \(f\in S_k\setminus\{0\}\)，还得到 \(|f(x+\mathrm{i}y)|\le C e^{-2\pi m y}\)、\(m\ge1\)，说明尖点处的消失在上半平面中表现为指数衰减。

由局部分解立即可得，在各函数均非恒零且在尖点亚纯时，
\[
 v_\infty(fg)=v_\infty(f)+v_\infty(g),\quad
 v_\infty(f/g)=v_\infty(f)-v_\infty(g).
\]
若 \(f+g\not\equiv0\)，则 \(v_\infty(f+g)\ge\min(v_\infty(f),v_\infty(g))\)；两阶相等时首项可能相消而使不等式严格。后面从 \(E_4^3\) 与 \(E_6^2\) 相减构造判别式，正是有意制造这样的相消。

最后，\(q=0\) 的阶数是尖点坐标中的阶数，不是把 \(\tau=\infty\) 当作普通有限点后读出的次数。有限的椭圆固定点另有 \(u^2\)、\(u^3\) 商坐标，它们与这里的指数坐标分工不同。\secref{b3:sec:2305}将把这些局部阶数放入同一条零点计数公式，而每一项所用的坐标都会保持明确。

% !TeX root = ../../Book3.tex
\section{\texorpdfstring{\(E_4,E_6\)}{E₄, E₆} 与判别式}
\label{b3:sec:2304}

已经知道 \(G_4,G_6\) 是模形式，但它们的常数项分别含有 \(\zeta(4),\zeta(6)\)。把常数项规范为一，可以更直接地比较系数。随后，\(E_4^3\) 与 \(E_6^2\) 具有相同的权和常数项，二者相减便产生一个尖点模形式。这个构造还将恢复前章三次曲线的判别式，因而能把尖点处的消失与上半平面内部的非退化性联系起来。

\subsection{归一化 Eisenstein 级数}
\label{b3:subsec:230401}

\begin{definition}[归一化 Eisenstein 级数]\label{b3:def:23-normalized-eisenstein}
设 \(k\ge2\) 为整数，\(\tau\in\mathfrak H\)，且
\[
 G_{2k}(\tau)=\sum_{(m,n)\in\mathbb Z^2\setminus\{(0,0)\}}(m\tau+n)^{-2k},
 \quad \zeta(2k)=\sum_{n=1}^{\infty}n^{-2k}.
\]
定义归一化的 Eisenstein 级数
\[
 E_{2k}(\tau)\coloneqq\frac{G_{2k}(\tau)}{2\zeta(2k)}.
\]
\end{definition}
\symidx{\(E_{2k}(\tau)\)：常数项为一的 Eisenstein 级数}

分母是正实数，所以这个定义不会除以零。\cref{b3:prop:23-eisenstein-transformation}与\cref{b3:prop:23-eisenstein-fourier}说明 \(E_{2k}\in M_{2k}\)，并给出
\begin{equation}\label{b3:eq:23-normalized-eisenstein-fourier}
 E_{2k}(\tau)=1+
 \frac{(2\pi\mathrm{i})^{2k}}{(2k-1)!\,\zeta(2k)}
 \sum_{n=1}^{\infty}\sigma_{2k-1}(n)q^n,
 \quad q=e^{2\pi\mathrm{i}\tau}.
\end{equation}
它的常数项等于一，故不是尖点模形式。要把权四、六的系数写成具体整数，还需要算出相应的偶数幂和。这个计算只用本卷已建立的余切展开。

\begin{proposition}[三个偶数幂和的值]\label{b3:prop:23-zeta-even-values}
对实数 \(s>1\)，令 \(\zeta(s)=\sum_{n=1}^{\infty}n^{-s}\)。则
\[
 \zeta(2)=\frac{\pi^2}{6},\quad
 \zeta(4)=\frac{\pi^4}{90},\quad
 \zeta(6)=\frac{\pi^6}{945}.
\]
\end{proposition}
\begin{proof}
由\cref{b3:prop:20-sine-product}，在 \(0<|z|<1\) 内，
\[
 \pi\cot\pi z
 =\frac1z+\sum_{n=1}^{\infty}\frac{2z}{z^2-n^2}.
\]
对每个正整数 \(n\) 展开几何级数，
\[
 \frac{2z}{z^2-n^2}
 =-2\sum_{r=1}^{\infty}\frac{z^{2r-1}}{n^{2r}}.
\]
若 \(|z|\le\rho<1\)，右侧双和的绝对值之和至多为
\[
 2\sum_{n=1}^{\infty}\frac{\rho/n^2}{1-\rho^2/n^2}
 \le\frac{2\rho}{1-\rho^2}\sum_{n=1}^{\infty}n^{-2}<\infty.
\]
所以可以交换两次求和，得到
\begin{equation}\label{b3:eq:23-cotangent-zeta}
 \pi\cot\pi z=\frac1z-2\sum_{r=1}^{\infty}\zeta(2r)z^{2r-1},
 \quad 0<|z|<1.
\end{equation}

现在用正弦、余弦的 Taylor 级数独立计算左端。置 \(t=\pi z\)，并写
\[
 t\cot t=1+At^2+Bt^4+Ct^6+O(t^8).
\]
因为 \(t\cot t=\cos t/(\sin t/t)\)，所以
\[
 \begin{aligned}
 &\left(1-\frac{t^2}{6}+\frac{t^4}{120}
              -\frac{t^6}{5040}+O(t^8)\right)
   \left(1+At^2+Bt^4+Ct^6+O(t^8)\right)\\
 &\hspace{2em}=1-\frac{t^2}{2}+\frac{t^4}{24}
                    -\frac{t^6}{720}+O(t^8).
 \end{aligned}
\]
逐次比较二、四、六次项，得到
\[
 A-\frac16=-\frac12,\quad
 B-\frac A6+\frac1{120}=\frac1{24},\quad
 C-\frac B6+\frac A{120}-\frac1{5040}=-\frac1{720}.
\]
第一式给 \(A=-1/3\)；代入第二式得 \(B=-1/45\)，再代入第三式得 \(C=-2/945\)。换回 \(z\)，
\[
 \pi\cot\pi z=\frac1z-\frac{\pi^2}{3}z
                    -\frac{\pi^4}{45}z^3
                    -\frac{2\pi^6}{945}z^5+O(z^7).
\]
与\eqref{b3:eq:23-cotangent-zeta}比较对应系数，就得到三个数值。
\end{proof}

这里先由正规收敛的部分分式得到幂和系数，再由初等 Taylor 除法算出同一函数的系数。整个证明不需要关于 \(\zeta(s)\) 的函数方程；后面学习解析数论时，可以把这几个数值作为已知的正实轴数据使用。

\subsection{\texorpdfstring{\(E_4,E_6\)}{E4,E6}}
\label{b3:subsec:230402}

\begin{proposition}[权四与权六的 Eisenstein 级数]\label{b3:prop:23-e4-e6}
设 \(\tau\in\mathfrak H\)，\(q=e^{2\pi\mathrm{i}\tau}\)，并按\cref{b3:def:23-normalized-eisenstein}定义 \(E_4,E_6\)。则 \(E_4\in M_4\)、\(E_6\in M_6\)，且
\begin{equation}\label{b3:eq:23-e4-e6-fourier}
 \begin{aligned}
 E_4(\tau)&=1+240\sum_{n=1}^{\infty}\sigma_3(n)q^n,\\
 E_6(\tau)&=1-504\sum_{n=1}^{\infty}\sigma_5(n)q^n.
 \end{aligned}
\end{equation}
这些幂级数在 \(|q|<1\) 内局部正规收敛，并有
\begin{equation}\label{b3:eq:23-e4-e6-initial}
 \begin{aligned}
 E_4&=1+240q+2160q^2+6720q^3+O(q^4),\\
 E_6&=1-504q-16632q^2-122976q^3+O(q^4).
 \end{aligned}
\end{equation}
特别地，二者都非恒零，且尖点消失阶均为零。
\end{proposition}
\begin{proof}
由\eqref{b3:eq:23-normalized-eisenstein-fourier}，权四的系数因子为
\[
 \frac{(2\pi\mathrm{i})^4}{3!\,\zeta(4)}
 =\frac{16\pi^4}{6\cdot(\pi^4/90)}=240,
\]
而权六的系数因子为
\[
 \frac{(2\pi\mathrm{i})^6}{5!\,\zeta(6)}
 =\frac{-64\pi^6}{120\cdot(\pi^6/945)}=-504.
\]
负号来自 \(\mathrm{i}^6=-1\)。一般 Fourier 展开已经保证了变换律与尖点全纯性，因此它们确实属于所述模形式空间。

前三个正整数的约数给出
\[
 \begin{aligned}
 &\sigma_3(1)=1,\quad \sigma_3(2)=1+2^3=9,\quad
   \sigma_3(3)=1+3^3=28,\\
 &\sigma_5(1)=1,\quad \sigma_5(2)=1+2^5=33,\quad
   \sigma_5(3)=1+3^5=244.
 \end{aligned}
\]
代入即可得到首项展开。常数项均为一，故二者不是零函数，也不在尖点消失。
\end{proof}

这些整数系数来自约数和及刚刚计算的常数。一般模形式可以乘任意复常数，所以“是模形式”本身没有强迫系数为整数。\(\tau\) 的函数恒等式也不能仅凭几个初项相同就判定；\secref{b3:sec:2305}的价公式将给出可以比较有限多个系数的精确依据。

\subsection{判别式模形式}
\label{b3:subsec:230403}

寻找尖点模形式时，可以先让常数项相消。但相减的两个模形式还必须具有相同的权，否则变换因子无法共同提出。权四与权六的最小正公倍数为十二，因此首先比较 \(E_4^3\) 与 \(E_6^2\)。二者常数项都是一，一次项却分别是 \(720q\) 与 \(-1008q\)，所以差的首项为 \(1728q\)。除以 \(1728\) 就使第一个非零系数也规范为一。

\begin{definition}[判别式模形式]\label{b3:def:23-delta}
设 \(\tau\in\mathfrak H\)，\(E_4,E_6\) 为\cref{b3:def:23-normalized-eisenstein}中常数项为一的权四、权六 Eisenstein 级数。函数
\[
 \Delta(\tau)\coloneqq
 \frac{E_4(\tau)^3-E_6(\tau)^2}{1728}
\]
称为\term[panbieshi-mo-xingshi]{判别式模形式}[discriminant modular form]。
\end{definition}
\symidx{\(\Delta(\tau)\)：判别式模形式}

\begin{proposition}[判别式的权与尖点阶数]\label{b3:prop:23-delta-cusp}
对上半平面上的归一化 Eisenstein 级数 \(E_4,E_6\)，令 \(\Delta=(E_4^3-E_6^2)/1728\)。则 \(\Delta\) 是全模群的非零权十二尖点模形式，并且
\begin{equation}\label{b3:eq:23-delta-fourier}
 \Delta(\tau)=q-24q^2+252q^3+O(q^4),
 \quad q=e^{2\pi\mathrm{i}\tau}.
\end{equation}
因此 \(v_\infty(\Delta)=1\)。
\end{proposition}
\begin{proof}
乘积的权相加，所以 \(E_4^3,E_6^2\in M_{12}\)，其差也属于 \(M_{12}\)。两者的常数项相等，故差的常数项为零，得到 \(\Delta\in S_{12}\)。

为确定消失阶，不能停在常数项抵消；还需计算下一项。由\eqref{b3:eq:23-e4-e6-initial}，
\[
 \begin{aligned}
 E_4^3&=1+720q+179280q^2+16954560q^3+O(q^4),\\
 E_6^2&=1-1008q+220752q^2+16519104q^3+O(q^4).
 \end{aligned}
\]
其中平方、立方的二次项分别来自
\[
 2(-16632)+(-504)^2=220752,\quad
 3\cdot2160+3\cdot240^2=179280;
\]
三次项分别来自
\[
 \begin{aligned}
 &2(-122976)+2(-504)(-16632)=16519104,\\
 &3\cdot6720+6\cdot240\cdot2160+240^3=16954560.
 \end{aligned}
\]
相减除以 \(1728\)，就得到\eqref{b3:eq:23-delta-fourier}。一次项系数为一，说明函数非恒零，而且在局部参数 \(q\) 中恰有一阶零点。
\end{proof}

消失阶是在 \(q=0\) 的局部坐标中计数，不是在变量 \(\tau\) 的某个有限点求 Taylor 阶数。具体地，\(\Delta/q=1-24q+252q^2+O(q^3)\) 在 \(q=0\) 附近全纯且非零，因此
\[
 \Delta(x+\mathrm{i}y)
 =e^{2\pi\mathrm{i}x}e^{-2\pi y}
   \left(1+O(e^{-2\pi y})\right)
 \quad(y\longrightarrow\infty),
\]
误差对 \(x\in\mathbb R\) 一致。这证明它在足够高的半平面内没有零点；要排除其余内部零点，还需另一个理由。

\subsection{判别式的两种归一化}
\label{b3:subsec:230404}

前章使用的是格的判别式 \(D(\Lambda)=g_2(\Lambda)^3-27g_3(\Lambda)^2\)，本节使用的是常数已经规范化的模形式 \(\Delta(\tau)\)。两者都具有十二次尺度关系，但数值并不相同。先把比例算准，前章的非退化结论才可以直接用于这里。

\begin{proposition}[格判别式与模判别式]\label{b3:prop:23-delta-nonvanishing}
设 \(\tau\in\mathfrak H\)，\(\Lambda_\tau=\mathbb Z+\mathbb Z\tau\)，并定义
\[
 g_2(\Lambda_\tau)=60\sum_{0\ne\omega\in\Lambda_\tau}\omega^{-4},
 \quad
 g_3(\Lambda_\tau)=140\sum_{0\ne\omega\in\Lambda_\tau}\omega^{-6}.
\]
令 \(D(\Lambda_\tau)=g_2(\Lambda_\tau)^3-27g_3(\Lambda_\tau)^2\)，且
\(\Delta(\tau)=(E_4(\tau)^3-E_6(\tau)^2)/1728\)，其中 \(E_4,E_6\) 为常数项一的 Eisenstein 级数。则
\begin{equation}\label{b3:eq:23-discriminant-normalization}
 D(\Lambda_\tau)=(2\pi)^{12}\Delta(\tau)\ne0.
\end{equation}
更一般地，对每个 \(a\in\mathbb C\setminus\{0\}\)，有
\[
 D(a\Lambda_\tau)=a^{-12}(2\pi)^{12}\Delta(\tau).
\]
\end{proposition}
\begin{proof}
由\eqref{b3:eq:23-g-invariants}及偶数幂和的值，
\[
 g_2(\Lambda_\tau)
 =120\zeta(4)E_4=\frac{4\pi^4}{3}E_4,\quad
 g_3(\Lambda_\tau)
 =280\zeta(6)E_6=\frac{8\pi^6}{27}E_6.
\]
因此
\[
 \begin{aligned}
 D(\Lambda_\tau)
 &=\frac{64\pi^{12}}{27}\left(E_4^3-E_6^2\right)\\
 &=\frac{64\cdot1728}{27}\pi^{12}\Delta
  =4096\pi^{12}\Delta=(2\pi)^{12}\Delta.
 \end{aligned}
\]
因为 \(\operatorname{Im}\tau>0\)，\(\Lambda_\tau\) 是真正的秩二格。\cref{b3:prop:22-discriminant-scaling}已证明，对每一个这样的格，三个非零半周期的 \(\wp\) 值 \(e_1,e_2,e_3\) 互异，而且
\[
 D(\Lambda_\tau)=16\prod_{1\le j<\ell\le3}(e_j-e_\ell)^2\ne0.
\]
于是 \(\Delta(\tau)\ne0\)。同一命题的缩放律 \(D(a\Lambda)=a^{-12}D(\Lambda)\) 给出最后一式。
\end{proof}

这里的“判别式”沿用前章的 \(D\) 约定。若使用多项式 \(4X^3-g_2X-g_3\) 的通常多项式判别式，数值还要乘以 \(16\)。固定非零因子不影响是否有重根，却影响规范化后的精确等式。

\ProseParagraphBreak

尖点处消失与内部没有零点并不冲突。沿 \(\tau=\mathrm{i}y\)、\(y\to\infty\)，\(q\to0\)，由 Fourier 展开，
\[
 g_2\longrightarrow\frac{4\pi^4}{3},\quad
 g_3\longrightarrow\frac{8\pi^6}{27},\quad D\longrightarrow0.
\]
极限的三次多项式为
\[
 4X^3-\frac{4\pi^4}{3}X-\frac{8\pi^6}{27}
 =4\left(X-\frac{2\pi^2}{3}\right)
    \left(X+\frac{\pi^2}{3}\right)^2.
\]
它确实出现了重根，但这是趋向尖点的极限，不是某个有限 \(\tau\in\mathfrak H\) 所对应的格曲线。

\ProseParagraphBreak

至此，\(\Delta\) 在上半平面内部处处可除，而在尖点恰消去一次 \(q\) 因子。这两个性质将在\secref{b3:sec:2305}一起使用：把尖点模形式除以 \(\Delta\) 时，内部不会产生极点，尖点处能否全纯则由消失阶决定。

% !TeX root = ../../Book3.tex
\section{价公式与模形式环}
\label{b3:sec:2305}

普通辐角原理在基本区域上仍有效，但圆弧配对时函数出现权因子，椭圆固定点又只贡献部分圆周。价公式就是把这两种修正逐项记录的结果。

\subsection{基本域中的零点}
\label{b3:subsec:230501}

对非恒零亚纯模形式，记 \(v_p(f)\) 为上半平面局部参数中的零点阶数，极点取负数。模群变换中的因子在上半平面非零，故阶数只依赖轨道。取每条普通轨道一个代表；\(\mathrm{i}\) 与 \(\rho\) 单列。尖点阶数按 \(q\) 参数定义。

\ProseParagraphBreak

固定点的阶数还有同余限制。例如 \(f(-1/\tau)=\tau^kf(\tau)\) 在 \(\mathrm{i}\) 附近比较首项，得到 \((-1)^{v_i}=\mathrm{i}^k\)。权为二时必在 \(\mathrm{i}\) 消失，足以排除非零权二全纯形式；一般同余关系也可以精化维数计数。

\subsection{价公式}
\label{b3:subsec:230502}

\begin{theorem}[价公式]\label{b3:thm:23-valence}
若 \(f\) 是非零权 \(k\) 亚纯模形式，在尖点也亚纯，则
\[
v_\infty(f)+\frac12v_{\mathrm{i}}(f)+\frac13v_\rho(f)
+\sum_{p\ne\mathrm{i},\rho}v_p(f)=\frac{k}{12},
\]
其中求和按普通模群轨道计数。
\end{theorem}
\begin{proof}
先假定边界没有零极点。在高度 \(Y\) 截断基本域，取正向边界：底部圆弧从 \(\rho\) 经 \(\mathrm{i}\) 走到 \(\rho+1\)，顶部由右向左。周期性使竖边的 \(f'/f\) 积分相消；顶部积分趋于 \(-2\pi\mathrm{i}v_\infty(f)\)。
把底弧分为左半弧 \(L:\rho\to\mathrm{i}\) 与右半弧 \(R:\mathrm{i}\to\rho+1\)。映射 \(S\) 将 \(L\) 映成反向的 \(R\)。微分变换律给
\(\mathrm{d}\log f(S\tau)=\mathrm{d}\log f(\tau)+k\,\mathrm{d}\tau/\tau\)。故两弧积分之和等于 \(-k\int_L\mathrm{d}\tau/\tau\)。左弧的辐角从 \(2\pi/3\) 减至 \(\pi/2\)，结果为 \(k\pi\mathrm{i}/6\)。辐角原理除以 \(2\pi\mathrm{i}\) 得内部零极点总数等于 \(k/12-v_\infty\)。
若边界有零极点，先删去其小圆邻域，并将成对边界上的删去部分配成映射下的对应段。\(f'/f=v_p/(\tau-p)+O(1)\)，所以小弧积分的极限等于其角度占 \(2\pi\) 的比例乘 \(2\pi\mathrm{i}v_p\)。普通边界点的两份半圆合成整圆；\(\mathrm{i}\) 只占角度 \(\pi\)；\(\rho,\rho+1\) 两角合计 \(2\pi/3\)。于是分别得到系数一、二分之一、三分之一。被删除的小段对权修正积分的贡献趋零，顶部仍用同一 \(q\) 展开估计。取小圆半径趋零即得所示公式。
\end{proof}
这一结果称为\term[jia-gongshi]{价公式}[valence formula]。Zagier 的《Elliptic Modular Forms and Their Applications》\cite[Proposition 2, pp. 9--10]{Zagier2008Elliptic}可用于核对上述固定点权重；本书在围道计算中显式保留了边界方向。

\subsection{维数论证}
\label{b3:subsec:230503}

\begin{theorem}\label{b3:thm:23-modular-ring}
全模群的分次模形式环为 \(M_* =\mathbb C[E_4,E_6]\)，两个生成元代数独立。对偶数 \(k\ge0\)，
\[
\dim M_k=\begin{cases}\lfloor k/12\rfloor,&k\equiv2\pmod{12},\\
\lfloor k/12\rfloor+1,&k\not\equiv2\pmod{12}.
\end{cases}
\]
负权及奇权空间均为零。
\end{theorem}
\begin{proof}
价公式各项非负，故负权无非零形式；奇权已由 \(-I\) 排除。权零形式无零极点，或用紧化基本域上的最大值原理，可知为常数。权二时 \(v_i\ge1\)，其贡献至少 \(1/2>2/12\)，故空间为零。
因 \(\Delta\) 内部无零点且尖点一阶消失，除以 \(\Delta\) 给出 \(S_k=\Delta M_{k-12}\)。每个偶数 \(k\ge4\) 都可写成 \(4a+6b\)，其中 \(a,b\ge0\)：\(k\equiv0\pmod4\) 时取 \(b=0\)，\(k\equiv2\pmod4\) 时取 \(b=1\)。从 \(f\in M_k\) 减去其常数项乘 \(E_4^aE_6^b\)，余项属于 \(S_k\)。对权归纳，得到所有形式均为两生成元的多项式。每次常数项映射满射到 \(\mathbb C\)，故 \(\dim M_k=1+\dim M_{k-12}\)；结合 \(0,2,4,6,8,10\) 六个初值即得公式。
为证代数独立，先把任何多项式关系按权分解。对固定 \(\tau\)，代入无限多个矩阵 \(\left(\begin{smallmatrix}1&0\\n&1\end{smallmatrix}\right)\)，得到不同权部分乘 \((n\tau+1)^k\) 的多项式恒等式；不同次数的系数必须各为零。因此只需排除同权关系。同权单项式约去一个共同因子后，成为 \(E_4^3/E_6^2\) 的多项式关系。这个亚纯函数非常数，因为两者常数项均一而一次系数分别为 \(720,-1008\)。非常值亚纯函数在没有极点的某小圆盘上像含开集，不可能满足非零常系数多项式，故关系为零。
\end{proof}
例如 \(M_8\) 一维，常数项为一的 \(E_8\) 必等于 \(E_4^2\)。比较 Fourier 系数即得到约数和卷积恒等式。有限维数把一个无穷级数恒等式缩减为有限个系数的核对；减去已匹配项后，价公式保证余项必须为零。

\ProseParagraphBreak

模形式环还允许构造一个保持模性的微分运算。普通导数的变换律含有额外项，必须先计算这项，而不能直接把导数称为高二权形式。记 \(D=(2\pi\mathrm{i})^{-1}\mathrm{d}/\mathrm{d}\tau=q\mathrm{d}/\mathrm{d}q\)，并令
\[
 P(\tau)=\frac{D\Delta(\tau)}{\Delta(\tau)}.
\]
判别式在上半平面无零点，故 \(P\) 全纯。尖点处 \(\Delta=q(1-24q+\cdots)\)，所以 \(P=1-24q+O(q^2)\)。对 \(\Delta(\gamma\tau)=(c\tau+d)^{12}\Delta(\tau)\) 求对数导数，链式法则给
\[
 P(\gamma\tau)=(c\tau+d)^2P(\tau)
       +\frac{12c(c\tau+d)}{2\pi\mathrm{i}}.
\]
额外的加项说明它不是权二模形式；这与 \(M_2=0\) 完全一致。它的任务是抵消其他函数导数中的同类加项。
\ProseParagraphBreak
对权 \(k\) 形式 \(f\) 求导，得到
\[
 (Df)(\gamma\tau)=(c\tau+d)^{k+2}Df(\tau)
       +\frac{kc}{2\pi\mathrm{i}}(c\tau+d)^{k+1}f(\tau).
\]
故 \(D_kf=Df-(k/12)Pf\) 满足权 \(k+2\) 的变换律。其 \(q\) 展开没有负幂；若 \(f\) 是尖点形式，两个乘积也仍无常数项。因此 \(D_k\) 把 \(M_k\) 映入 \(M_{k+2}\)，并保持尖点子空间。这是对普通导数的一项具体修正，全部常数都由权十二的判别式决定。
\ProseParagraphBreak
现在用已经证明的低权维数即可得微分恒等式。\(D_4E_4\in M_6\)，其常数项为 \(-1/3\)，所以 \(D_4E_4=-E_6/3\)。同理 \(D_6E_6\in M_8=\mathbb C E_4^2\)，常数项为 \(-1/2\)。于是
\[
 DE_4=\frac{PE_4-E_6}{3},\qquad
 DE_6=\frac{PE_6-E_4^2}{2}.
\]
再对 \(P\) 的变换律求导，可得 \(DP-P^2/12\) 为权四模形式。具体而言，令 \(J=c\tau+d\)，导数变换中的额外两项是 \(2cJ^3P/(2\pi\mathrm{i})\) 与 \(12c^2J^2/(2\pi\mathrm{i})^2\)；展开 \(P(\gamma\tau)^2/12\) 恰给相同两项，因而完全相消。其常数项为 \(-1/12\)，故
\[
 DP=\frac{P^2-E_4}{12}.
\]
这些式子展示有限维性怎样将微分信息转为代数关系。它们不依赖尚未证明的无限乘积公式，也没有从前几项相同就直接宣称两函数相同；关键是先证明差属于一个已知的一维模形式空间。
\ProseParagraphBreak
也可反过来从这组微分方程递推系数。若写 \(P=1+\sum p_nq^n\)、\(E_4=1+\sum a_nq^n\)、\(E_6=1+\sum b_nq^n\)，每个 \(D\) 把第 \(n\) 项乘以 \(n\)，乘积的第 \(n\) 项则是有限卷积。给定前面的系数后，方程约束后续系数，而 Eisenstein 级数的显式约数和满足这些约束。这是可计算的一致性检验，但如果要据此从零定义三个函数，仍须另外证明所得形式级数的收敛；本章已有全纯构造，故无需倒置这一步。

% !TeX root = ../../Book3.tex
\section{模函数与 \texorpdfstring{\(j\)}{j}-不变量}
\label{b3:sec:2306}

\(E_4\) 与 \(E_6\) 随格基的改变而乘上不同因子，不能直接用它们的数值比较两个复环面。若先把权调成相同，再取商，因子便会约去。问题随之变为：能否找到一个这样的商，使它的值不仅在等价的格上相同，而且足以辨认全部格比例类？下面将用价公式回答这个逆向问题，再给所得参数空间建立具体的复坐标。

\subsection{模函数}
\label{b3:subsec:230601}

\begin{definition}\label{b3:def:23-modular-function}
设 \(\mathfrak H=\{\tau\in\mathbb C:\operatorname{Im}\tau>0\}\)，群 \(\mathrm{SL}_2(\mathbb Z)\) 按分式线性变换作用，\(h\) 为 \(\mathfrak H\) 上的亚纯函数。若
\[
 h\left(\frac{a\tau+b}{c\tau+d}\right)=h(\tau)
 \quad\left(\begin{pmatrix}a&b\\c&d\end{pmatrix}\in\mathrm{SL}_2(\mathbb Z)\right),
\]
且周期函数 \(h\) 的 \(q=e^{2\pi\mathrm i\tau}\) 表达式在 \(q=0\) 亚纯，则称 \(h\) 为全模群的\term[mo-hanshu]{模函数}[modular function]。这里的亚纯函数不包括恒取 \(\infty\) 的映射，允许恒零函数。
\end{definition}

尖点条件允许有限个负次幂；它不要求在尖点有界。若 \(h\not\equiv0\)，则在充分高的半平面有
\[
 h(\tau)=q^m\bigl(a_m+a_{m+1}q+\cdots\bigr),
 \qquad m\in\mathbb Z,\quad a_m\ne0.
\]
正的 \(m\) 表示尖点零点，负的 \(m\) 表示尖点极点。由于 \(|q|=e^{-2\pi\operatorname{Im}\tau}\)，一个 \(r\) 阶尖点极点沿高半平面的增长量级为 \(e^{2\pi r\operatorname{Im}\tau}\)。因此“在无穷远亚纯”说的是指数坐标中的有限阶行为，不能照搬平面坐标 \(1/\tau\) 下的极点定义。

\ProseParagraphBreak
若 \(f,g\in M_k\) 且 \(g\not\equiv0\)，则 \(f/g\) 为模函数。上半平面中的变换因子相消；在任一有限点把分子、分母各自提出其零点幂次，商只有有限阶零点或极点。尖点处同样对两个收敛的 \(q\) 幂级数作局部分解。这说明商在有限点和尖点都亚纯，不能仅凭形式上“权相减为零”省略这两处检查。

\ProseParagraphBreak
非零模函数的倒数、两个模函数的和与积仍为模函数，因而它们组成域。本章稍后将把这个域确定为 \(\mathbb C(j)\)。这里 \(\mathbb C(j)\) 指常系数有理函数代入 \(j\) 所得的函数域，而 \(M_*\) 是按权分次的全纯模形式环；前者允许零权函数有极点，后者用非零权容纳丰富的全纯对象，两者承担不同职责。

\subsection{\texorpdfstring{\(j\)}{j}-不变量}
\label{b3:subsec:230602}

判别式 \(\Delta\) 在上半平面没有零点，且在尖点恰好一阶消失。因此用它作分母，可以把全部极点集中到唯一的尖点。分子须有权十二；选 \(E_4^3\) 还使首项系数为一，这固定了一个方便的归一化。

\begin{definition}\label{b3:def:23-j-invariant}
对 \(\tau\in\mathfrak H\)，令 \(E_4,E_6\) 为全模群上常数项归一为一的权四、权六 Eisenstein 级数，并令 \(\Delta=(E_4^3-E_6^2)/1728\)。定义
\term[j-bubianliang]{\(j\)-不变量}[j-invariant]为
\[
 j(\tau)=\frac{E_4(\tau)^3}{\Delta(\tau)}.
\]
\end{definition}
\symidx{\(j(\tau)\)：椭圆模不变量}

两个权十二因子相除，故 \(j\) 对全模群不变。由前节的 Fourier 展开，
\[
 E_4^3=1+720q+179280q^2+O(q^3),\qquad
 \Delta=q\bigl(1-24q+252q^2+O(q^3)\bigr).
\]
先求第二个括号的倒数：若写成 \(1+b_1q+b_2q^2+O(q^3)\)，乘回原式比较系数得到 \(b_1=24\)、\(b_2=24^2-252=324\)。再与分子相乘，得到
\begin{equation}\label{b3:eq:23-j-fourier}
 j=q^{-1}+744+196884q+O(q^2).
\end{equation}
这里的常数项为 \(720+24\)，一次项为 \(179280+720\cdot24+324\)。归一化中的 \(1728\) 使尖点主部恰为 \(q^{-1}\)；若改取 \(j/1728\)，分类性质不变，但特殊值和全部展开系数都须相应改变。

\Needspace{6\baselineskip}
\begin{theorem}[\texorpdfstring{\(j\)}{j} 对模群轨道的分离]\label{b3:thm:23-j-orbits}
设 \(j=E_4^3/\Delta\)，其中 \(E_4,E_6\) 为全模群归一化 Eisenstein 级数，\(\Delta=(E_4^3-E_6^2)/1728\)，并令 \(\rho=e^{2\pi\mathrm i/3}\)。对每个 \(w\in\mathbb C\)，集合
\(\{\tau\in\mathfrak H:j(\tau)=w\}\) 恰为一条 \(\mathrm{PSL}_2(\mathbb Z)\) 轨道。若 \(p\) 在这条轨道上，\(j-w\) 在上半平面参数 \(\tau-p\) 中的零点阶数为
\[
 \begin{cases}
 3,&p\text{ 与 }\rho\text{ 模群等价},\\
 2,&p\text{ 与 }\mathrm i\text{ 模群等价},\\
 1,&p\text{ 为其他点}.
 \end{cases}
\]
此外 \(j(\rho)=0\)、\(j(\mathrm i)=1728\)。
\end{theorem}
\begin{proof}
固定 \(w\)，考察
\[
 F_w=E_4^3-w\Delta\in M_{12}.
\]
它的常数项为一，故非恒零，且 \(v_\infty(F_w)=0\)。由于 \(\Delta\) 内部无零，\(F_w\) 与 \(j-w\) 在上半平面具有相同零点及阶数。\cref{b3:thm:23-valence}于是给出
\[
 \frac12v_{\mathrm i}(F_w)+\frac13v_\rho(F_w)
 +\sum_{p\ne\mathrm i,\rho}v_p(F_w)=1,
\]
普通点仍按轨道取一个代表。

这个等式本身只给加权总数，必须再知道每条非空零点轨道至少贡献一，才能推出唯一性。普通点的零点阶数为正整数，已有这一性质。在椭圆固定点使用 \(j-w\) 的不变性：由\cref{b3:prop:23-elliptic-orders}的权零情形，它在 \(\mathrm i\) 的阶数为二的倍数，在 \(\rho\) 的阶数为三的倍数。因此这两个位置除以稳定子阶数后的贡献也为正整数。

总数为一首先排除了零点集合为空，又只允许一条零点轨道，并强迫其商参数中的阶数为一。换回上半平面参数，就得到所述三种阶数。最后，由\cref{b3:cor:23-eisenstein-zeros}，\(E_4(\rho)=0\)，而 \(E_6(\mathrm i)=0\)。代入
\[
 j=\frac{E_4^3}{\Delta},\qquad
 j-1728=\frac{E_6^2}{\Delta}
\]
便得到两个特殊值。
\end{proof}

\ProseParagraphBreak
此证明同时给出存在和唯一性。不能先假设某个 \(w\) 被取到，再只证明原像唯一；价公式右侧的正数保证每个有限 \(w\) 都确实有原像。它也说明基本域边界上的两个代表不应算成两个复结构：总数始终按群轨道计，而非按闭基本域中互异点的个数计。

\ProseParagraphBreak
例如 \(E_4^3-1728\Delta=E_6^2\) 的零点全部位于 \(\mathrm i\) 的轨道，每处为二阶；\(E_4^3\) 的零点位于 \(\rho\) 的轨道，每处为三阶。若 \(w\notin\{0,1728\}\)，则 \(E_4^3-w\Delta\) 只有一条普通零点轨道，每处为简单零点。三个情形的加权总数都是一。

\subsection{复环面的 \texorpdfstring{\(j\)}{j} 分类}
\label{b3:subsec:230603}

前面的轨道结论是在参数 \(\tau\) 上得到的。要把它解释成复环面的分类，还须排除一种可能：两个格不成比例，却有某种非线性的全纯坐标变化把它们的商曲面同构起来。平面覆盖使这一问题可以转为整函数的导数估计。

\Needspace{6\baselineskip}
\begin{proposition}[复环面之间的全纯映射]\label{b3:prop:23-torus-affine}
设 \(\Lambda,\Lambda'\subset\mathbb C\) 为秩二格，\(T_\Lambda=\mathbb C/\Lambda\)、\(T_{\Lambda'}=\mathbb C/\Lambda'\) 带有平面商坐标所定义的复结构。每个全纯映射 \(F:T_\Lambda\to T_{\Lambda'}\) 都可写为
\[
 F([z]_\Lambda)=[az+b]_{\Lambda'},\qquad
 a,b\in\mathbb C,\quad a\Lambda\subseteq\Lambda'.
\]
反之，每个满足上述包含条件的仿射式都定义全纯映射。映射 \(F\) 双全纯当且仅当 \(a\ne0\) 且 \(a\Lambda=\Lambda'\)；若还要求 \(F([0])=[0]\)，则可取 \(b=0\)。
\end{proposition}
\begin{proof}
记商投影为 \(\pi_\Lambda,\pi_{\Lambda'}\)。\cref{b3:prop:22-torus-charts}已证明这两个投影是全纯覆盖：取足够小的圆盘，其不同格平移互不相交，投影在每一片上双全纯。平面单连通，故\cref{b3:thm:11-map-lifting}应用于 \(F\circ\pi_\Lambda:\mathbb C\to T_{\Lambda'}\)，给出连续提升 \(\widetilde F:\mathbb C\to\mathbb C\)。提升在局部等于 \(\pi_{\Lambda'}\) 的全纯逆分支与 \(F\circ\pi_\Lambda\) 的复合，因此为整函数。

对每个 \(\lambda\in\Lambda\)，两个点 \(z\) 与 \(z+\lambda\) 的投影相同，所以
\[
 \widetilde F(z+\lambda)-\widetilde F(z)\in\Lambda'.
\]
左端随 \(z\) 连续，右端是离散集合；平面连通使这个差对 \(z\) 为常数。对 \(z\) 求导，便知 \(\widetilde F'\) 为整的 \(\Lambda\)-周期函数。它在闭基本平行四边形上有界，再由周期性在整个平面有界，Liouville 定理使 \(\widetilde F'=a\) 为常数。积分得到 \(\widetilde F(z)=az+b\)，原来的差式给 \(a\lambda\in\Lambda'\)。

反向的良定义也正由这个条件保证：改变代表 \(z\mapsto z+\lambda\)，像代表只增加 \(a\lambda\in\Lambda'\)；局部商坐标中是仿射全纯函数。

若 \(F\) 双全纯，便不恒定，所以 \(a\ne0\)。若有 \(\lambda'\in\Lambda'\setminus a\Lambda\)，则 \([\lambda'/a]_\Lambda\ne[0]_\Lambda\)，却与 \([0]_\Lambda\) 映到同一点，违背单射。因此 \(a\Lambda=\Lambda'\)。反之，等式成立时 \(w\mapsto(w-b)/a\) 良定义地下降为逆映射，故双全纯。若 \(F([0])=[0]\)，则 \(b\in\Lambda'\)，在商中可以去掉这一项。
\end{proof}

这里并未预先假设 \(F\) 保持加法。结论说明，先用目标环面的平移把 \(F([0])\) 移回原点，所得全纯映射自动由复线性映射下降，因此自动是群同态。若只要求实线性同构，则任意两组实基都能相互映到，全部秩二格的商在拓扑上都是环面；复线性这一条件才保留共形结构。

\Needspace{6\baselineskip}
\begin{theorem}[复环面的模参数分类]\label{b3:thm:23-torus-j-classification}
设 \(\Lambda=\omega_1(\mathbb Z+\mathbb Z\tau)\)、\(\Lambda'=\omega_1'(\mathbb Z+\mathbb Z\tau')\)，其中 \(\omega_1,\omega_1'\ne0\)，\(\tau,\tau'\in\mathfrak H\)。以下条件等价：
\begin{enumerate}
\item 复环面 \(\mathbb C/\Lambda\) 与 \(\mathbb C/\Lambda'\) 双全纯；
\item 存在 \(a\in\mathbb C^*\) 使 \(a\Lambda=\Lambda'\)；
\item \(\tau'\) 与 \(\tau\) 属于同一 \(\mathrm{PSL}_2(\mathbb Z)\) 轨道；
\item \(j(\tau')=j(\tau)\)。
\end{enumerate}
每个有限复数均为某个复环面的 \(j\) 值。
\end{theorem}
\begin{proof}
前两项等价由\cref{b3:prop:23-torus-affine}给出，其中常数平移不改变格条件。\cref{b3:prop:23-lattice-orbits}已用正向整数基变换证明：归一化格 \(\mathbb Z+\mathbb Z\tau\) 与 \(\mathbb Z+\mathbb Z\tau'\) 成复比例，当且仅当两参数由行列式一的整数矩阵联系；缩放因子 \(\omega_1,\omega_1'\) 不影响这一判别。最后两项等价以及有限值的满射性由\cref{b3:thm:23-j-orbits}得到。
\end{proof}

\ProseParagraphBreak
对格 \(\Lambda_\tau\)，\secref{b3:sec:2304}给出的常数换算使
\[
 j(\tau)=1728\frac{g_2(\Lambda_\tau)^3}
 {g_2(\Lambda_\tau)^3-27g_3(\Lambda_\tau)^2}.
\]
分子、分母在格缩放下同乘十二次逆幂，故也可直接用任意未归一化的格计算这个比值。方格的 \(g_3=0\) 给 \(j=1728\)，六角格的 \(g_2=0\) 给 \(j=0\)。特殊值表示额外的旋转对称，但两类环面仍各自只有一个同构类，分类定理没有在这些位置失效。

\ProseParagraphBreak
若 \(Y,Y'>1\) 且 \(Y\ne Y'\)，则 \(\mathrm iY,\mathrm iY'\) 是基本域内部不同点，因此对应的矩形环面不双全纯，\(j(\mathrm iY)\ne j(\mathrm iY')\)。另一方面，\(\mathrm iY\) 与 \(\mathrm i/Y\) 由 \(S\) 联系，代表同一环面。这个例子说明“参数不同”与“复结构不同”之间还隔着换基。

\ProseParagraphBreak
由\cref{b3:thm:22-uniformization}，每个格环面已经与它的 Weierstrass 光滑射影三次曲线双全纯对应，原点对应无穷远点。因此上述结论也分类这些由格构造的曲线的保原点解析同构。\(j\) 相等意味着相应复结构同构，并不要求两条仿射方程的系数逐项相等。

\subsection{模曲线 \texorpdfstring{\(X(1)\)}{X(1)} 的几何解释}
\label{b3:subsec:230604}

目前 \(j\) 已给出轨道集合与 \(\mathbb C\) 的双射。集合双射还没有说明连续性、局部复坐标和无穷远点。为理解“模参数空间也是一张 Riemann 曲面”，下面把这些结构逐一建立。

令 \(G=\mathrm{PSL}_2(\mathbb Z)\)，\(Y(1)=G\backslash\mathfrak H\) 带商拓扑，\(\pi:\mathfrak H\to Y(1)\) 为投影。\secref{b3:sec:2301}已证明紧集有限相交性质及局部商模型。普通点附近可取投影单射的小圆盘；在 \(p=\mathrm i\) 或 \(\rho\) 附近，先用
\[
 u_p(\tau)=\frac{\tau-p}{\tau-\overline p}
\]
把稳定子作用化为旋转，再用 \(u_p^e\) 作商坐标，分别取 \(e=2,3\)。其他与这两点等价的位置先经群元搬回即可。这是\cref{b3:prop:23-stabilizer-chart}的具体内容；不能在固定点直接把 \(\pi\) 当成无分歧覆盖。

\ProseParagraphBreak
这些图之间的转移是全纯的。在没有固定点的小片上，选定局部提升后，转移由 Möbius 变换及全纯坐标复合给出；椭圆点中心与普通小片不会同时相交，在去心重叠处可选 \(u_p^e\) 的一个根分支。同一椭圆轨道可固定一个上述坐标图；缩小此图或改用等价代表时，群变换与稳定子旋转相容，相应不变全纯函数按幂级数下降到商坐标。因此这些局部图给出相容的复结构。

\Needspace{6\baselineskip}
\begin{proposition}[唯一尖点的紧化]\label{b3:prop:23-modular-compactification}
设 \(G=\mathrm{PSL}_2(\mathbb Z)\) 按分式线性变换作用于 \(\mathfrak H\)。在 \(Y(1)=G\backslash\mathfrak H\) 中加入一个点 \(\infty\)，规定它的邻域基为
\[
 U_Y=\{\infty\}\cup\pi\{\tau\in\mathfrak H:\operatorname{Im}\tau>Y\},
 \qquad Y>1.
\]
用 \(q=e^{2\pi\mathrm i\tau}\) 将 \(U_Y\) 识别为 \(|q|<e^{-2\pi Y}\)，并以 \(q=0\) 对应新增点。它与普通点坐标、椭圆点坐标 \(u_p^e\) 共同定义一个紧连通 Riemann 曲面，记为 \(X(1)\)。
\end{proposition}
\begin{proof}
\cref{b3:prop:23-cusp-chart}保证当 \(Y>1\) 时，高半平面内只有整数平移可能识别两个点。因此 \(q\) 在 \(U_Y\setminus\{\infty\}\) 上良定义且为一一的局部全纯坐标，像恰为去心圆盘。补上零点就得到完整圆盘。在两个尖点图的重叠上坐标相同；与内部图的重叠上，\(q\) 的导数 \(2\pi\mathrm i q\) 非零，局部对数给出全纯逆，故转移相容。

先证明内部轨道商是 Hausdorff 的。取不属于同一轨道的 \(p,q\in\mathfrak H\)，并取它们的紧邻域 \(K,L\)。由\cref{b3:prop:23-proper-action}，满足 \(\gamma K\cap L\ne\varnothing\) 的群元只有有限个。每个这样的群元均有 \(\gamma p\ne q\)，故可缩小 \(p,q\) 的开邻域，使其在这个群元下不相交；有限次取交后，得到 \(U\subset K,V\subset L\)，使所有这些群元均满足 \(\gamma U\cap V=\varnothing\)。候选集合以外的群元本来就不可能相交。因此 \(\pi(U),\pi(V)\) 是不交的开邻域：商投影为开映射，因为一个开集全部群平移的并仍开。

还须把新尖点与任一内部点分开。固定 \(\tau_0\in\mathfrak H\)，选一个含它的小紧邻域 \(K\)，其中 \(\eta\le\operatorname{Im}\tau\le B\)、\(\eta>0\)。对 \(\tau=x+\mathrm i y\in K\) 及任意群元，若 \(c\ne0\)，则
\[
 \operatorname{Im}(\gamma\tau)
 =\frac{y}{(cx+d)^2+c^2y^2}\le\frac1{c^2y}\le\eta^{-1};
\]
若 \(c=0\)，则虚部仍为 \(y\le B\)。取 \(Y>\max\{1,B,\eta^{-1}\}\)，\(\pi(K)\) 便与 \(U_Y\) 不交，从而尖点与内部点也有不交邻域。

将上半平面的可数基用开商投影送到 \(Y(1)\)，可得 \(Y(1)\) 的可数基。再添上整数高度的 \(U_Y\)，得到 \(X(1)\) 的可数基。因此上述局部图给出 Hausdorff、第二可数的一维复流形。\(Y(1)\) 是连通上半平面的连续像，而每个尖点邻域都与它相交，故添点后仍连通。

最后取 \(Y>1\)。截断基本域
\[
 K_Y=\{\tau\in\mathfrak H:|\operatorname{Re}\tau|\le1/2,
 \ |\tau|\ge1,\ \operatorname{Im}\tau\le Y\}
\]
是紧集：它闭有界，且虚部至少为 \(\sqrt3/2\)。它的投影覆盖高度不超过 \(Y\) 的基本域部分。其余部分加上尖点，在一个稍大尖点图中恰由闭圆盘 \(|q|\le e^{-2\pi Y}\) 参数化，也是紧的。这两个紧集的并覆盖 \(X(1)\)，故 \(X(1)\) 紧致。
\end{proof}

这个紧曲面称为全模群的\term[mo-quxian]{模曲线}[modular curve] \(X(1)\)。此处只讨论这一特定群的解析商及其紧化，不需要一般模空间的代数几何定义。单独加入的尖点不对应一个秩二格环面；当归一化基本域中的虚部趋向无穷时，它记录的是离开所有紧参数集的极限端。

\ProseParagraphBreak
模函数恰好下降为 \(X(1)\) 上的亚纯函数。普通点的下降来自群不变性；在阶数为 \(e\) 的椭圆点，局部 Laurent 级数在 \(u\mapsto\zeta_eu\) 下不变，故只含 \(u^{en}\)，改写成 \(u^e\) 的 Laurent 级数后仍只有有限个负幂。尖点处的亚纯性正是定义中的 \(q\) 条件。反过来，曲面上的亚纯函数拉回 \(\mathfrak H\) 后不变，并在尖点给出同样的有限负幂展开。这一步还说明：有限点在上半平面中的阶数，要除以稳定子阶数，才是商曲面上的阶数。

\Needspace{6\baselineskip}
\begin{theorem}[模曲线的球面坐标]\label{b3:thm:23-j-sphere}
令 \(X(1)\) 为全模群的上述紧化商曲面，\(j=E_4^3/\Delta\)，其中 \(E_4,E_6\) 为全模群归一化 Eisenstein 级数，\(\Delta=(E_4^3-E_6^2)/1728\)。令 \(j(\infty)=\infty\)，则 \(j\) 诱导双全纯映射
\[
 j:X(1)\longrightarrow\widehat{\mathbb C\kern0pt}.
\]
因此全模群的模函数域为 \(\mathbb C(j)\)。其中在整个上半平面全纯的模函数恰为 \(j\) 的多项式，若在尖点也全纯，则必为常数。
\end{theorem}
\begin{proof}
在内部，\cref{b3:thm:23-j-orbits}证明每个有限值恰对应一个商点，并在商坐标中为简单取值。换言之，\(j-j(p)\) 在普通坐标或 \(u_p^e\) 中都有非零一次系数，局部逆函数定理给全纯逆。尖点处，由\eqref{b3:eq:23-j-fourier}，
\[
 \frac1j=q\bigl(1-744q+O(q^2)\bigr),
\]
一次系数也非零；目标无穷远坐标是 \(1/j\)，故此处同样有全纯局部逆。内部没有极点，尖点是无穷远值的唯一原像。因此全局双射及各处全纯局部逆拼成所述双全纯映射。

若 \(h\) 是模函数，它下降到 \(X(1)\) 后与 \(j^{-1}\) 复合，是球面上的亚纯函数。\cref{b3:thm:03-sphere-rational}说明这种函数必为有理函数：球面紧致使极点只有有限多个，减去有限点的主要部分后，只剩一个在无穷远至多有极点的整函数，即多项式。因此 \(h=R(j)\)，其中 \(R\in\mathbb C(t)\)。反向把任意有理函数代入这个球面坐标，当然给出模函数。

若 \(h\) 在上半平面无极点，\(R\) 就不能有有限极点，因为每个有限值都被 \(j\) 取到；有理函数没有有限极点，当且仅当它是多项式。若它还在尖点无极点，该多项式在无穷远也全纯，只能为常数。
\end{proof}

\ProseParagraphBreak
例如一个模函数若只在尖点有至多二阶极点，就具有 \(Aj^2+Bj+C\) 的形式。\(j^2=q^{-2}+1488q^{-1}+\cdots\)，因此先用 \(j^2\) 匹配二阶主部，再用 \(j\) 匹配一阶主部，剩余仅是常数。这样的有限消元由函数域定理保证，不是仅凭若干 Fourier 系数作猜测。

\ProseParagraphBreak
尖点亚纯性在这里不可删除。函数 \(e^j\) 在上半平面全纯且不变，但在尖点有本性奇点：因为 \(1/j\) 是该点的一个全纯局部坐标，\(e^j\) 在此坐标中就是 \(e^{1/z}\)。它既不是模函数，也不是 \(j\) 的有理函数。另一方面，\(j\) 自己已经说明“内部全纯且不变”不强迫常值；必须再控制唯一缺失的尖点。

\ProseParagraphBreak
\(j\) 在上半平面有二阶、三阶局部分歧，与它在 \(X(1)\) 上双全纯并不冲突。先商掉稳定子之后，平方或立方参数已经吸收这些重数，\(j\) 在商坐标中只剩一阶变化。至此，格的换基、模形式的权、价公式中的分数权重和复环面的一个复数参数，便由明确的局部计算衔接起来。

% 本章小结（不计入编号）
\begin{chaptersummary}
模形式同时约束上半平面的对称性与尖点增长。价公式把权转成零点总数，进而给出有限维性、\(E_4,E_6\) 的生成性与 \(\Delta\) 的最小消失阶。权零的比值 \(j\) 消去格尺度，将紧化参数空间识别为球面；这些有限维与系数比较的方法将在平方和问题中继续出现。
\end{chaptersummary}


\BookExercises

% 八题自拟：变换律、固定点、有限维性及系数比较的应用。
\begin{exercise}\label{b3:ex:23-1}
求 \(\mathcal F\) 对面积元 \(y^{-2}\,\mathrm{d}x\,\mathrm{d}y\) 的面积。
\end{exercise}
\begin{solution}
先积 \(y\)，得 \(\int_{-1/2}^{1/2}(1-x^2)^{-1/2}\,\mathrm{d}x=\pi/3\)。无穷远端有有限面积，但并非通常 Euclid 面积有限。
\end{solution}
\begin{exercise}\label{b3:ex:23-2}
确定 \(E_4\) 与 \(E_6\) 的全部零点轨道及其局部阶数。
\end{exercise}
\begin{solution}
\(ST\) 固定 \(\rho\)，权四变换因子不为一，故 \(E_4(\rho)=0\)。价公式总数 \(1/3\) 强迫该点为单零，别处无零。\(S\) 在 \(\mathrm{i}\) 的权六因子为 \(-1\)，故 \(E_6\) 在该点消失；总数 \(1/2\) 强迫单零且无其他轨道。
\end{solution}
\begin{exercise}\label{b3:ex:23-3}
证明两个 \(M_k\) 中的形式，只要从常数项起前 \(\lfloor k/12\rfloor+1\) 个系数相同，就完全相同。
\end{exercise}
\begin{solution}
差的尖点阶数大于 \(k/12\)。若差非零，价公式左侧已经仅由这一项超过右侧，矛盾。此为充分界，在 \(k\equiv2\pmod{12}\) 等情形可以改进。
\end{solution}
\begin{exercise}\label{b3:ex:23-4}
求 \(M_{14}\)、\(S_{24}\) 的基。
\end{exercise}
\begin{solution}
维数分别为一、二。前者由常数项一的 \(E_4^2E_6\) 张成。后者为 \(\Delta M_{12}\)，故可取 \(\Delta E_4^3,\Delta^2\)；两者尖点阶分别为一、二，线性独立。
\end{solution}
\begin{exercise}\label{b3:ex:23-5}
用 \(E_4^2=E_8\) 推导约数和卷积恒等式。
\end{exercise}
\begin{solution}
由一般 Eisenstein 展开有 \(E_8=1+480\sum\sigma_7(n)q^n\)，其中常数由 \(\zeta(8)=\pi^8/9450\) 得到，此值也由余切展开比较七次项算出。比较 \(q^n\) 得
\(\sum_{m=1}^{n-1}\sigma_3(m)\sigma_3(n-m)=(\sigma_7(n)-\sigma_3(n))/120\)。空和 \(n=1\) 时为零，右侧也零。
\end{solution}
\begin{exercise}\label{b3:ex:23-6}
证明在上半平面全纯且在尖点有界的模函数必常数，并给出删去尖点条件的反例。
\end{exercise}
\begin{solution}
有界性使 \(q=0\) 可去，函数下降为紧曲面 \(X(1)\) 上的全纯函数，最大值原理使其常数。\(j\) 在上半平面全纯且不变，但在尖点有极点，所以提供反例。
\end{solution}
\begin{exercise}\label{b3:ex:23-7}
若模函数只有尖点处至多 \(n\) 阶极点，证明它是次数至多 \(n\) 的 \(j\) 多项式。
\end{exercise}
\begin{solution}
利用 \(j=q^{-1}+744+\cdots\)，逐次减去 \(j^n,j^{n-1},\ldots,j\) 的适当倍数可消去全部负幂。余下函数在紧曲面全纯，故常数。最高极点阶数同时给多项式次数的上界。
\end{solution}
\begin{exercise}\label{b3:ex:23-8}
证明 \(j-1728=E_6^2/\Delta\)，并确定 \(j\) 在 \(\rho\)、\(j-1728\) 在 \(\mathrm{i}\) 的消失阶。
\end{exercise}
\begin{solution}
第一式由 \(E_4^3-E_6^2=1728\Delta\) 得到。判别式内部无零，且前题给 \(E_4,E_6\) 各自的单零，故两消失阶分别为三、二。在商空间的立方、平方局部参数中都成为一阶，这与 \(j\) 双全纯并不矛盾。
\end{solution}
